\documentclass[11pt]{article}
\usepackage[margin=1.1in]{geometry}
\usepackage{amsmath,amssymb,amsthm}
\usepackage{enumitem}
\usepackage{booktabs}
\usepackage{xcolor}
\definecolor{linkblue}{RGB}{25,60,120}
\usepackage[colorlinks=true,linkcolor=linkblue,citecolor=linkblue,urlcolor=linkblue]{hyperref}
\usepackage[capitalize]{cleveref}
\emergencystretch=2em

\newtheorem{theorem}{Theorem}[section]
\newtheorem{lemma}[theorem]{Lemma}
\newtheorem{proposition}[theorem]{Proposition}
\newtheorem{corollary}[theorem]{Corollary}
\newtheorem{fact}[theorem]{Fact}
\theoremstyle{definition}
\newtheorem{definition}[theorem]{Definition}
\newtheorem{problem}[theorem]{Open Problem}
\theoremstyle{remark}
\newtheorem{remark}[theorem]{Remark}
\crefname{problem}{Problem}{Problems}
\Crefname{problem}{Problem}{Problems}
\crefname{lemma}{Lemma}{Lemmas}
\Crefname{lemma}{Lemma}{Lemmas}
\crefname{proposition}{Proposition}{Propositions}
\Crefname{proposition}{Proposition}{Propositions}
\crefname{corollary}{Corollary}{Corollaries}
\Crefname{corollary}{Corollary}{Corollaries}
\crefname{definition}{Definition}{Definitions}
\Crefname{definition}{Definition}{Definitions}
\crefname{remark}{Remark}{Remarks}
\Crefname{remark}{Remark}{Remarks}

\newcommand{\F}{\mathbb{F}}
\newcommand{\B}{B}
\newcommand{\RS}{\mathrm{RS}}
\newcommand{\eca}{\varepsilon_{\mathrm{ca}}}
\newcommand{\emca}{\varepsilon_{\mathrm{mca}}}
\newcommand{\dmin}{\delta_{\min}}
\newcommand{\dist}{\Delta}
\newcommand{\distS}{\Delta_S}
\newcommand{\Lst}{\Lambda}
\newcommand{\Syn}{\operatorname{Syn}}
\newcommand{\Hb}{\mathrm{H}_2}
\newcommand{\eone}{e_1}

\title{A Universal Field-Size Cap and Prize-Facing Completion Program\\ for Mutual Correlated Agreement on Smooth Reed--Solomon Domains\\[2pt]\large Version 10: quotient, extension-pole, and Hankel rank-drop certificates}
\author{Przemek Chojecki\\ \small ulam.ai}
\date{July 1, 2026\\{\small v10 integrated milestone upgrade}}

\begin{document}
\maketitle

\begin{abstract}
We prove a field-size-universal upper bound on the proximity parameter achievable in the grand mutual-correlated-agreement (MCA) challenge of Arnon--Boneh--Fenzi \cite{ABF26}. Let $\F$ be any finite field with $|\F|<2^{256}$, let $D$ be a smooth multiplicative evaluation domain of order $n$ contained in a subfield $\B\subseteq\F$ (with $\B=\F$ allowed), and let $C=\RS[\F,D,k]$ have $k=\rho n\le2^{40}$ and $\rho\in\{1/2,1/4,1/8,1/16\}$. If $2^{10}\mid n$ for $\rho\in\{1/2,1/4,1/8\}$, and $2^{11}\mid n$ for $\rho=1/16$, then
\[
\emca(C,\delta)
\;>\;\frac{1}{2k}\Bigl(1-\frac n{|\F|}\Bigr)
\;\ge\;2^{-86}\;\gg\;2^{-128}
\]
for every $\delta\in[1-\rho-2^{-9},1-\rho)$ at rates $1/2,1/4,1/8$, and for every $\delta\in[1-\rho-2^{-10},1-\rho)$ at rate $1/16$. The same lower bound holds for $\eca(C,\delta)$ on the integer-admissible subinterval ending at $1-\rho-1/n$. If $|\F|\ge2n$, the displayed lower bound improves to $2^{-42}$. Consequently $\delta^*_C(2^{-128})\le 1-\rho-2^{-9}$ for $\rho\in\{1/2,1/4,1/8\}$ and $\delta^*_C(2^{-128})\le 1-\rho-2^{-10}$ for $\rho=1/16$, throughout the $|\F|<2^{256}$ challenge envelope. This is a prize-facing negative result: it does not determine the exact threshold, but it rules out an entire near-capacity band for every challenge-envelope row satisfying the printed divisor hypothesis.

The mechanism composes a pigeonhole lower bound on quotient-locator list fibers---a slack-two and subfield-pigeonhole sharpening of \cite[Main Thm.\,(d)]{Cho26a}---with a self-contained deep-point list-to-correlated-agreement conversion. This conversion is the same shallow consequence previously imported from Crites--Stewart \cite[Thm.~2]{CS25}, but it is proved here directly with the integer-radius condition. The composition itself was first carried out, conditionally and with weaker constants, in the universal-cap theorem of \cite{Cho26b} (error $2^{-42}$ at gap $2^{-11}$, assuming $|\F|\ge2n$); the present note sharpens the gap, removes the field-size condition, and extends the cap to extension fields. No value-set or equidistribution input is used, so the per-fiber collision problem isolated in \cite{Cho26b} is routed around rather than solved. Version 10 adds a quantitative inverse form of the deep-point conversion, heaviest quotient-remainder prefix fibers, an extension-pole conversion producing genuinely extension-valued witness lines, an exact divisor-block coefficient formula for the support-union ledger, a gcd/lcm quotient-image ledger that coalesces duplicate finite, projective, and curve parameters inside declared quotient branches, and a canonical maximal-minor gcd ledger for regular overdetermined Hankel buckets. We instantiate the bound at the KoalaBear-sextic parameters of \cite[\S6.3]{ABF26}, including the interleaved rows of their Tables~2--3, give an independent slacked variant through \cite[Thm.~1.9]{BCHKS25}, and prove a two-radius refinement of the subfield-confinement theorem of \cite{Cho26b}, showing that over extension fields every explicit base-rational locator line is inert while the extension-pole floor gives a concrete source of genuinely $\F$-valued certifying lines whenever a large list is present.
\end{abstract}

\section{Introduction}\label{sec:intro}

Proximity testing for Reed--Solomon codes underlies the soundness of modern SNARKs, and the sharpest available soundness accounting is phrased through \emph{mutual correlated agreement} (MCA): a random linear combination $f_1+\gamma f_2$ being $\delta$-close to the code should force $f_1,f_2$ to be simultaneously explained by codewords on a \emph{common} agreement set, except with small probability $\emca(C,\delta)$ over $\gamma$. The survey of Arnon, Boneh and Fenzi \cite{ABF26} crystallizes the state of the art into a \emph{grand MCA challenge}: for $C=\RS[\F,L,k]$ over a smooth domain $L$, with rate $\rho\in\{1/2,1/4,1/8,1/16\}$, degree $k\le2^{40}$, and field size $|\F|<2^{256}$, determine the largest proximity parameter $\delta^*$ at which $\emca(C,\delta)\le2^{-128}$ can hold.

Positive results reach the Johnson bound $1-\sqrt\rho$. On the negative side, \cite{Cho26b} showed $\delta^*_C(2^{-128})<1-\rho-1/64$ for all prime fields of size up to roughly $2^{150}$ (per-rate reaches between $2^{150.6}$ and $2^{161.8}$), by exhibiting quotient-locator witnesses whose bad-slope sets are value sets of restricted symmetric sums over small subgroups; it further proved a \emph{conditional} universal cap over every field of size at most $2^{256}$ --- error $2^{-42}$ at gap $2^{-11}$, assuming $|\F|\ge2n$ --- by composing quotient-core lists with the same conversion isolated and proved directly here (its universal-cap theorem). Above the value-set reach, what remained open was the \emph{error-one} question, identified there with the \emph{per-fiber collision problem}: controlling, at a \emph{fixed} large prime, the collisions of characteristic-zero locator values under reduction. Extension fields, the setting actually deployed (\cite[\S6.3]{ABF26} works over a sextic extension of KoalaBear), posed an additional obstruction: by the subfield-confinement theorem of \cite{Cho26b}, refined in \cref{sec:subfield}, every explicit witness of \cite{Cho26a,Cho26b} is $\B$-rational and contributes only $p^{-5}<2^{-154}$ over the deployed extension.

This note sharpens the universal cap --- gap $2^{-9}$ (resp.\ $2^{-10}$) in place of $2^{-11}$, error $2^{-86}$ with no condition relating $n$ and $|\F|$, and extension fields handled directly --- for every field below $2^{256}$ at once, by changing what is counted.


\begin{theorem}[informal; see \cref{thm:main,cor:grand}]\label{thm:informal}
Let $\F$ be any finite field with $|\F|<2^{256}$, let $D\subseteq\F^\times$ be a smooth multiplicative evaluation domain of order $n$, and let $C=\RS[\F,D,k]$ have $k=\rho n\le2^{40}$ with $\rho\in\{1/2,1/4,1/8,1/16\}$. Assume $2^{10}\mid n$ for $\rho\in\{1/2,1/4,1/8\}$, and $2^{11}\mid n$ for $\rho=1/16$. Then
\[
\emca(C,\delta)>2^{-86}\gg2^{-128}
\]
for every $\delta\in[1-\rho-2^{-9},1-\rho)$ at rates $1/2,1/4,1/8$, and for every $\delta\in[1-\rho-2^{-10},1-\rho)$ at rate $1/16$. The same lower bound holds for $\eca(C,\delta)$ on the integer-admissible subinterval ending at $1-\rho-1/n$. If $|\F|\ge2n$, the lower bound is $\ge2^{-42}$. Hence $\delta^*_C(2^{-128})\le1-\rho-2^{-9}$ for $\rho\in\{1/2,1/4,1/8\}$ and $\delta^*_C(2^{-128})\le1-\rho-2^{-10}$ for $\rho=1/16$.
\end{theorem}

The proof composes two ingredients. The first (\cref{lem:fiber}) is a pigeonhole lower bound on list sizes at a single received word: the quotient locators of \cite{Cho26a} attach to each $\ell$-subset $A$ of the order-$N$ quotient coset $Q=D^{a}$ a codeword $c_A$ at distance $\le1-\rho-t/N$ from the word $u_{z_A}=x^{k+ta}+z_Ax^{k+(t-1)a}$, where $z_A=-\eone(A)$; pigeonholing the $\binom N\ell$ subsets over their $\le|\B|$ slope values produces one word $u_z$ carrying a list of size $\binom N\ell/|\B|$. Two refinements of \cite[Main Thm.\,(d)]{Cho26a} matter here: we run the construction at \emph{slack two} ($t=2$), so that the listed codewords have degree $\le k$ and live in $\RS[\F,D,k+1]$ without any divisibility demand on $k+1$ (which is odd in all challenge instantiations, while $a$ is a $2$-power); and we pigeonhole over the \emph{subfield} $\B$ containing $D$ rather than over $\F$, which is what makes the extension-field case go through. The second ingredient is the direct deep-point conversion of \cref{thm:A}: in the integer-error range $f\le n-k-1$, if $\eca(C,\delta)$ is below $\eta(\frac1k-\frac n{k|\F|})$, then \emph{every} list of $\RS[\F,D,k+1]$ at radius $\delta$ has size at most $\lceil|\F|\,\eca/(1-\eta)\rceil$. Taking $\eta=\frac12$, the heavy fiber of \cref{lem:fiber} violates this whenever $\binom N\ell\ge|\B|(|\F|/k+1)$ --- an inequality with hundreds of bits of slack throughout the challenge envelope --- and the contrapositive yields $\eca(C,\delta)>\frac1{2k}(1-n/|\F|)$ throughout the integer-admissible range. The MCA cap then persists up to capacity by support-wise MCA monotonicity.

No value-set counting, exponential-sum estimate, or equidistribution hypothesis enters: pigeonhole fibers exist over every field unconditionally, and \cref{thm:A} converts list mass into correlated-agreement error directly. The cost of this route, relative to the error-one results of \cite{Cho26a,Cho26b} below $2^{150}$, is an error of order $1/k$ rather than $1-o(1)$, and a gap $2^{-9}$ at the first three prize rates (uniformly $2^{-10}$ across all four), rather than the fixed gap $1/64$; \cref{sec:discussion} quantifies the trade-off and what remains open.

Beyond the universal cap (\cref{cor:grand}, with a two-tier summary combining it with \cite{Cho26b}), we prove: a deployed-parameter corollary at the KoalaBear-sextic instantiation of \cite[\S6.3]{ABF26}, with $\emca>2^{-22}$ at gap $2^{-7}$ and the same $\eca$ lower bound on the integer-admissible subinterval (\cref{cor:deployed}), extended to every interleaved row of \cite[Tables~2--3]{ABF26} via an interleaving-transfer lemma (\cref{lem:inter,cor:rows}); an independent slacked variant through \cite[Thm.~1.9]{BCHKS25} for robustness (\cref{prop:slacked}); and a subfield-confinement lemma (\cref{lem:confine}) showing that for lines with $\B$-valued words over an extension $\F/\B$, every CA- or MCA-bad slope lies in $\B$. The latter has a structural consequence (\cref{cor:Fvalued}): at deployed parameters the lines certifying \cref{cor:deployed} are necessarily \emph{not} $\B$-rational, so the failure proved here is fiber-borne and currently non-constructive; exhibiting explicit certifying lines is posed as Open Problem~\ref{prob:explicit}.

\paragraph{Roadmap.} \Cref{sec:prelim} fixes definitions, aligned with \cite[Defs.~4.1, 4.3]{ABF26}, proves the direct deep-point conversion used for the main cap, and records the separate BCHKS/ABF slacked fallback. \Cref{sec:fiber} proves the fiber and interleaving lemmas and the conservative quotient-support upper ledger, \cref{sec:main} the main theorem, \cref{sec:grand} the challenge-level and deployed corollaries, and \cref{sec:subfield} subfield confinement. The new \cref{sec:prize-program} states exactly what this paper contributes toward the Proximity Prize \cite{ProximityPrize}, gives conditional completion theorems for the MCA and interleaved-list thresholds, and isolates the proof obligations that would turn the present cap into a full threshold determination. \Cref{sec:discussion} discusses scope, limits, calibration against the conjectural picture of \cite{Cho26b}, and open problems. Throughout, companion results of \cite{Cho26a,Cho26b} are cited by their printed theorem names rather than numbers, since the numbering of those manuscripts is not yet frozen.

\section{Preliminaries}\label{sec:prelim}

\paragraph{Codes and distances.} Throughout, $\B\subseteq\F$ are finite fields, $q:=|\F|$, and $D=\alpha H\subseteq\B^\times$ is a multiplicative coset of a cyclic subgroup $H$ of order $n$; we write $D$ for the evaluation domain (denoted $L$ in \cite{ABF26}). The subgroup case is $\alpha=1$, and taking $\B=\F$ is always allowed. For $k\le n$,
\[
\RS[\F,D,k]\;:=\;\bigl\{(p(x))_{x\in D}\;:\;p\in\F[X],\ \deg p<k\bigr\}\subseteq\F^n,
\qquad \rho:=k/n .
\]
For $u,v\in\F^n$, $\dist(u,v)$ is the relative Hamming distance, $\dist(u,C):=\min_{c\in C}\dist(u,c)$, and the relative minimum distance of $C=\RS[\F,D,k]$ is $\dmin(C)=1-\rho+1/n$. For $S\subseteq D$, $\distS(u,v)$ and $\distS(u,C)$ denote the corresponding restricted relative distances on $S$. Lists are denoted
\[
\Lst(C,\delta,u):=\{c\in C:\dist(u,c)\le\delta\},\qquad \Lst(C,\delta):=\max_{u\in\F^n}|\Lst(C,\delta,u)| .
\]
For $s\ge1$ the $s$-interleaved code is $C^{\equiv s}:=\{(c^{(1)},\dots,c^{(s)}):c^{(i)}\in C\}\subseteq(\F^n)^s$ with the column distance $\dist_s(F,G):=\frac1n\,\bigl|\{j:\exists i,\ f^{(i)}_j\ne g^{(i)}_j\}\bigr|$, and $\dist_s(F,C^{\equiv s}):=\min_{\mathbf c\in C^{\equiv s}}\dist_s(F,\mathbf c)$.

\begin{definition}[correlated agreement error; {\cite[Def.~4.1]{ABF26}} up to notation]\label{def:ca}
For $0\le\delta_{\mathrm{fld}}\le\delta_{\mathrm{int}}<1$,
\[
\eca(C,\delta_{\mathrm{fld}},\delta_{\mathrm{int}})\;:=\;\max_{f_1,f_2\in\F^n}\ \Pr_{\gamma\leftarrow\F}\Bigl[\ \dist(f_1+\gamma f_2,\,C)\le\delta_{\mathrm{fld}}\ \wedge\ \dist_2\bigl((f_1,f_2),\,C^{\equiv2}\bigr)>\delta_{\mathrm{int}}\ \Bigr],
\]
and the no-proximity-loss version is $\eca(C,\delta):=\eca(C,\delta,\delta)$. A slope $\gamma$ realizing the event for a fixed pair $(f_1,f_2)$ is called \emph{CA-bad} for that pair.
\end{definition}

\begin{definition}[mutual correlated agreement error, support-wise; {\cite[Def.~4.3]{ABF26}} up to notation]\label{def:mca}
For $\delta\in(0,1)$,
\[
\emca(C,\delta):=\max_{f_1,f_2\in\F^n}\Pr_{\gamma\leftarrow\F}\left[
\begin{array}{l}
\exists S\subseteq D,\ |S|\ge(1-\delta)n,\ \distS(f_1+\gamma f_2,C)=0,\\[2pt]
\text{and }\distS\bigl((f_1,f_2),C^{\equiv2}\bigr)>0
\end{array}
\right].
\]
Equivalently, $\gamma$ is MCA-bad for $(f_1,f_2)$ if some large agreement set $S$ explains the point $f_1+\gamma f_2$, but the same set $S$ does not simultaneously explain $f_1$ and $f_2$ by two codewords of $C$.
\end{definition}

\begin{definition}[challenge threshold]\label{def:dstar}
Following the grand MCA challenge of \cite[\S1]{ABF26}, for $\varepsilon^*\in(0,1)$ set
$\delta^*_C(\varepsilon^*):=\sup\{\delta\in(0,1-\rho):\emca(C,\delta)\le\varepsilon^*\}$, with $\sup\emptyset:=0$.
\end{definition}

\begin{fact}[{cf.\ \cite[Fact~4.5]{ABF26}}]\label{fact:chain}
For every $\delta\in(0,1)$: $\eca(C,\delta)\le\emca(C,\delta)$.
\end{fact}

\begin{proof}
Fix $(f_1,f_2)$ and let $\gamma$ be CA-bad: $\dist(f_1+\gamma f_2,C)\le\delta$ and $\dist_2((f_1,f_2),C^{\equiv2})>\delta$. Choose $S\subseteq D$ of size at least $(1-\delta)n$ on which $f_1+\gamma f_2$ agrees with a codeword of $C$. Since $(f_1,f_2)$ is not $\delta$-close to $C^{\equiv2}$, there is no set of size at least $(1-\delta)n$ on which $f_1$ and $f_2$ are simultaneously explained by codewords of $C$; in particular this fails on the chosen $S$. Thus $\gamma$ is MCA-bad for the same pair. Taking maxima over pairs gives the claim.
\end{proof}

\begin{lemma}[support-wise MCA monotonicity]\label{lem:mca-monotone}
For every linear code $C\subseteq\F^D$, the function
\[
        \delta\longmapsto \emca(C,\delta)
\]
is nondecreasing.
\end{lemma}

\begin{proof}
Fix a pair $(f_1,f_2)$ and suppose that $\gamma$ is MCA-bad at radius $\delta$. Then there is a set $S\subseteq D$ with
\[
        |S|\ge (1-\delta)n
\]
such that $f_1+\gamma f_2$ is explained by $C$ on $S$, but $(f_1,f_2)$ is not simultaneously explained by two codewords of $C$ on the same $S$. If $\delta'\ge\delta$, then
\[
        |S|\ge (1-\delta)n\ge (1-\delta')n,
\]
so the same set $S$ witnesses that $\gamma$ is MCA-bad at radius $\delta'$. Taking the maximum over pairs gives the claim.
\end{proof}

\paragraph{Conversion theorem.} The main composition below uses the deep-point list-to-CA conversion proved here directly. This is the same consequence previously imported from Crites--Stewart, with the integer-radius admissibility condition made explicit; see \cref{rem:import}. We keep the slacked BCHKS/ABF fallback as a separate imported route.

\begin{theorem}[deep-point list-to-CA conversion]
\label{thm:A}
Let $C=\RS[\F,D,k]$, let $C^+:=\RS[\F,D,k+1]$, let
$q:=|\F|$ and $n:=|D|$, and assume $q>n$.  Let $\delta\in(0,1)$ and set
$f_\delta:=\lfloor \delta n\rfloor$.  Assume
\[
        f_\delta \le n-k-1 .
\]
For any $\eta\in[0,1)$, if
\[
        \eca(C,\delta)
        \le
        \eta\Bigl(\frac1k-\frac n{kq}\Bigr),
\]
then
\[
        \Lst(C^+,\delta)
        \le
        \Bigl\lceil
        \frac{q\,\eca(C,\delta)}{1-\eta}
        \Bigr\rceil .
\]
\end{theorem}

\begin{proof}
Put $\epsilon:=\eca(C,\delta)$ and $f=f_\delta$.
Let $U:D\to\F$ be any received word and suppose that
$P_1,\ldots,P_L\in\F[X]_{<k+1}$ are distinct polynomials whose
restrictions to $D$ all lie in the Hamming ball of relative radius
$\delta$ around $U$.  Since distances are integral, each $P_i$ agrees
with $U$ on at least
\[
        a:=n-f
\]
points of $D$.  The hypothesis $f\le n-k-1$ gives $a\ge k+1$, in particular
$a>k$.

Let $\Omega:=\F\setminus D$, so $|\Omega|=q-n$.  For each
$\alpha\in\Omega$ define the simple-pole line
\[
        f_\alpha(x)=\frac{U(x)}{x-\alpha},
        \qquad
        g_\alpha(x)=-\frac1{x-\alpha}
        \qquad (x\in D).
\]
If $P_i$ agrees with $U$ on a set $S_i$ of size at least $a$, then for
$z_i(\alpha):=P_i(\alpha)$ and every $x\in S_i$,
\[
        f_\alpha(x)+z_i(\alpha)g_\alpha(x)
        =\frac{P_i(x)-P_i(\alpha)}{x-\alpha}.
\]
The right-hand side is the restriction of a polynomial of degree $<k$.
Thus each distinct value among the $P_i(\alpha)$ gives a slope for
which the line point is $\delta$-close to $C$.

We next check the correlated-agreement far condition.  If a polynomial
$G\in\F[X]_{<k}$ agreed with $g_\alpha$ on more than $k$ points of $D$,
then
\[
        (X-\alpha)G(X)+1
\]
would be a polynomial of degree at most $k$ with more than $k$ roots,
but its value at $X=\alpha$ is $1$, a contradiction.  Hence
$(f_\alpha,g_\alpha)$ is not $\delta$-close to $C^{\equiv2}$, because
any common explanation on a support of size at least $a>k$ would in
particular explain $g_\alpha$ there.  Therefore the distinct values of
$P_i(\alpha)$ are CA-bad slopes.

It remains to lower-bound the number of distinct values for some
$\alpha$.  For $i\ne j$, the polynomial $P_i-P_j$ is nonzero of degree
at most $k$, so it has at most $k$ roots in $\Omega$.  Let $M(\alpha)$
be the number of distinct values among $P_1(\alpha),\ldots,P_L(\alpha)$.
Over all $\alpha\in\Omega$, the total number of colliding unordered pairs
$(i,j)$ with $P_i(\alpha)=P_j(\alpha)$ is at most $k\binom L2$.
Equivalently, if the fiber sizes of the map $i\mapsto P_i(\alpha)$ are
$m_v(\alpha)$, then
\[
        \sum_{\alpha\in\Omega}\sum_v \binom{m_v(\alpha)}2
        \le k\binom L2 .
\]
Choose $\alpha\in\Omega$ for which the inner collision count is at most
$k\binom L2/(q-n)$.  If the fiber sizes of the map $i\mapsto P_i(\alpha)$ are
$m_1,\ldots,m_{M(\alpha)}$, then
\[
        \sum_r m_r=L,
        \qquad
        \sum_r m_r^2=L+2\sum_r\binom{m_r}{2}
        \le L+\frac{kL(L-1)}{q-n}.
\]
Cauchy--Schwarz therefore gives
\[
        L^2\le M(\alpha)\sum_r m_r^2,
        \qquad
        M(\alpha)\ge
        \frac{L(q-n)}{q-n+k(L-1)}
        \ge
        \frac{L(q-n)}{q-n+kL} .
\]
Consequently
\[
        \epsilon
        \ge
        \frac1q\cdot \frac{L(q-n)}{q-n+kL}.
\]
Solving this inequality for $L$ gives
\[
        L
        \le
        \frac{\epsilon q(q-n)}{q-n-k\epsilon q},
\]
provided $\epsilon<(q-n)/(kq)$.  Under the stated hypothesis,
$\epsilon\le\eta(q-n)/(kq)$ with $\eta<1$, so the denominator is at
least $(1-\eta)(q-n)$ and hence
\[
        L\le \frac{q\epsilon}{1-\eta}.
\]
Since $U$ and the list were arbitrary, the displayed ceiling bounds the
maximum list size of $C^+$ at radius $\delta$.
\end{proof}

\begin{remark}[relation to Crites--Stewart]
The displayed conversion is the only consequence of Crites--Stewart
needed by the universal-cap argument.  Earlier drafts imported it from
\cite[Thm.~2]{CS25}; the proof above is included to make the main cap
self-contained.  The separate slacked fallback through BCHKS/ABF remains
an independent imported route.
\end{remark}

\begin{proposition}[quantitative deep-list floor]
\label{prop:quantitative-deep-list-floor}
Let $C=\RS[\F,D,k]$, let $C^+:=\RS[\F,D,k+1]$, put
$q=|\F|$ and $n=|D|$, and assume $q>n$.  Let
$\delta_0\in[0,1-k/n)$, and suppose that for some received word
$U:D\to\F$ there are $L\ge1$ distinct codewords of $C^+$ in the
closed ball of radius $\delta_0$ around $U$.  Then, for every
$\delta\in[\delta_0,1-k/n)$,
\[
        \eca(C,\delta)
        \ge
        \frac{L(q-n)}{q(q-n+kL)} .
\]
For every such $\delta>0$,
\[
        \emca(C,\delta)
        \ge
        \frac{L(q-n)}{q(q-n+kL)}
\]
as well.  More precisely, some received line has at least
\[
        \left\lceil\frac{L(q-n)}{q-n+kL}\right\rceil
\]
CA-bad finite slopes at every such radius; for $\delta>0$, the same slopes are
MCA-bad.
\end{proposition}

\begin{proof}
Fix $\delta\in[\delta_0,1-k/n)$ and choose distinct
$P_1,\ldots,P_L\in\F[X]_{<k+1}$ in the $\delta_0$-ball around $U$.
They are also in the $\delta$-ball.  Since $\delta<1-k/n$, if
$f=\lfloor\delta n\rfloor$ then $f\le n-k-1$, so every codeword in this
$\delta$-ball agrees with $U$ on at least $a=n-f>k$ points.

Repeat the simple-pole construction in the proof of \cref{thm:A}.  For
$\alpha\in\Omega:=\F\setminus D$ set
\[
        f_\alpha(x)=\frac{U(x)}{x-\alpha},
        \qquad
        g_\alpha(x)=-\frac1{x-\alpha}
        \qquad (x\in D).
\]
For every distinct value of $P_i(\alpha)$, the word
$f_\alpha+P_i(\alpha)g_\alpha$ is $\delta$-close to $C$.  The pair
$(f_\alpha,g_\alpha)$ is not $\delta$-close to $C^{\equiv2}$, because a
common explanation on more than $k$ points would in particular explain
$g_\alpha$ there, and then $(X-\alpha)G(X)+1$ would be a polynomial of degree at most $k$ with more than $k$ roots and value $1$ at $\alpha$.
Thus the distinct values of the map $i\mapsto P_i(\alpha)$ are CA-bad
slopes.

For $i\ne j$, the polynomial $P_i-P_j$ has degree at most $k$ and is
nonzero, hence it has at most $k$ roots in $\Omega$.  Averaging over
$|\Omega|=q-n$ gives an $\alpha$ for which the number of colliding
unordered pairs among the $L$ values is at most $k\binom L2/(q-n)$.  If
$m_1,\ldots,m_M$ are the value multiplicities for this $\alpha$, then
\[
        \sum_r m_r=L,
        \qquad
        \sum_r m_r^2
        \le L+\frac{kL(L-1)}{q-n}.
\]
By Cauchy--Schwarz,
\[
        M\ge
        \frac{L^2}{L+kL(L-1)/(q-n)}
        =
        \frac{L(q-n)}{q-n+k(L-1)}
        \ge
        \frac{L(q-n)}{q-n+kL}.
\]
Since $M$ is an integer, $M\ge\lceil L(q-n)/(q-n+kL)\rceil$.  Dividing
by the $q$ possible finite slopes gives the CA lower bound.  For $\delta>0$, the MCA
bound follows from \cref{fact:chain}.
\end{proof}

\begin{corollary}[deep-list trigger and ceiling]
\label{cor:deep-list-trigger-ceiling}
In the setting of \cref{prop:quantitative-deep-list-floor}, define
\[
        \mathcal E_{q,k}(L)
        :=\frac{L(q-n)}{q(q-n+kL)},
        \qquad
        \mathcal N_{q,k}(L)
        :=\left\lceil\frac{L(q-n)}{q-n+kL}\right\rceil .
\]
Then $\mathcal E_{q,k}(L)$ is increasing in $L$, and
\[
        \mathcal E_{q,k}(L)>
        \frac1{2k}\left(1-\frac nq\right)
        \quad\Longleftrightarrow\quad
        L>\frac{q-n}{k}.
\]
Moreover
\[
        \mathcal E_{q,k}(L)<\frac1k\left(1-\frac nq\right)
\]
for every finite $L$, and the right-hand side is the limiting value as
$L\to\infty$.
\end{corollary}

\begin{proof}
The derivative of $L/(q-n+kL)$ is positive.  The trigger inequality is
\[
        \frac{L}{q-n+kL}>\frac1{2k},
\]
which is equivalent to $kL>q-n$.  The upper bound follows from
$q-n+kL>kL$.
\end{proof}

\begin{remark}[why this is stronger than the contrapositive]
\label{rem:quantitative-floor-vs-contrapositive}
\Cref{thm:A} is optimized for contradiction: once $L>(q-n)/k$, taking
$\eta=1/2$ already forces the printed $(2k)^{-1}(1-n/q)$ lower bound.
\Cref{prop:quantitative-deep-list-floor} keeps the whole interpolation count.
It is the form a staircase scanner should use when a quotient fiber is too small
to cross the $1/(2k)$ trigger but still contributes a nonzero explicit bad-slope
numerator, and also when a very large fiber improves the constant toward
$(1/k)(1-n/q)$.
\end{remark}

\begin{proposition}[punctured simple-pole deep-list floor]
\label{prop:punctured-simple-pole-floor}
Let $C=\RS[\F,D,k]$, let $C^+:=\RS[\F,D,k+1]$, put
$q=|\F|$ and $n=|D|$.  Let
\[
        \Omega_0\subseteq \F\setminus D,
        \qquad m:=|\Omega_0|>0 .
\]
Let $\delta_0\in[0,1-k/n)$, and suppose that for some received word
$U:D\to\F$ there are $L\ge1$ distinct codewords of $C^+$ in the closed ball of
radius $\delta_0$ around $U$.  Then, for every
$\delta\in[\delta_0,1-k/n)$,
\[
        \eca(C,\delta)
        \ge
        \frac{Lm}{q(m+kL)} .
\]
For every such $\delta>0$,
\[
        \emca(C,\delta)
        \ge
        \frac{Lm}{q(m+kL)} .
\]
More precisely, there is some $\alpha\in\Omega_0$ such that the simple-pole line
\[
        f_\alpha(x)=\frac{U(x)}{x-\alpha},
        \qquad
        g_\alpha(x)=-\frac1{x-\alpha}
        \qquad (x\in D)
\]
has at least
\[
        \left\lceil\frac{Lm}{m+kL}\right\rceil
\]
CA-bad finite slopes at every such radius; for $\delta>0$, the same slopes are
MCA-bad.
\end{proposition}

\begin{proof}
Repeat the proof of \cref{prop:quantitative-deep-list-floor}, but average only
over $\Omega_0$ instead of over all of $\F\setminus D$.  The simple-pole
argument works for every $\alpha\notin D$, so every distinct value among
$P_1(\alpha),\ldots,P_L(\alpha)$ gives a CA-bad slope.  For $i\ne j$, the
nonzero polynomial $P_i-P_j$ has degree at most $k$, hence has at most $k$ roots
inside the smaller set $\Omega_0$.  Therefore the total number of colliding
unordered pairs over all $\alpha\in\Omega_0$ is at most $k\binom L2$.  Choose
$\alpha\in\Omega_0$ for which the collision count is at most
$k\binom L2/m$.  If $M$ is the number of distinct values at this $\alpha$, the
same Cauchy--Schwarz calculation gives
\[
        M\ge
        \frac{Lm}{m+k(L-1)}
        \ge
        \frac{Lm}{m+kL} .
\]
Since $M$ is integral, this gives the displayed bad-slope numerator.  Dividing
by the $q$ finite slopes gives the CA lower bound, and the MCA statement for
$\delta>0$ follows from \cref{fact:chain}.
\end{proof}

\begin{corollary}[extension-pole deep-list floor]
\label{cor:extension-pole-deep-list-floor}
Assume in addition that $\B\subsetneq\F$, $D\subseteq\B$, and $|D|\ge2$.  Put
$b:=|\B|$.  In the setting of
\cref{prop:punctured-simple-pole-floor}, take
\[
        \Omega_0:=\F\setminus\B,
        \qquad m=q-b .
\]
Then for every $\delta\in[\delta_0,1-k/n)$,
\[
        \eca(C,\delta)
        \ge
        \mathcal E^{\rm ext}_{q,b,k}(L)
        :=
        \frac{L(q-b)}{q(q-b+kL)} .
\]
For $\delta>0$ the same bound holds for $\emca(C,\delta)$.  Moreover the
witnessing line may be chosen with pole $\alpha\in\F\setminus\B$, and its pole
vector $((x-\alpha)^{-1})_{x\in D}$ is not a scalar multiple over $\F$ of any
$\B$-valued vector.  In particular the witness line is genuinely
extension-valued, not merely a $\B$-rational line viewed over $\F$.
\end{corollary}

\begin{proof}
The lower bound is \cref{prop:punctured-simple-pole-floor} with
$\Omega_0=\F\setminus\B$.  It remains to prove genuine extension-valuedness.
Suppose that, for some $\lambda\in\F^\times$, all numbers
$\lambda/(x-\alpha)$ with $x\in D$ lay in $\B$.  Choose two distinct points
$x,y\in D$.  Then
\[
        \frac{y-\alpha}{x-\alpha}
        =
        \frac{\lambda/(x-\alpha)}{\lambda/(y-\alpha)}
        \in \B .
\]
Writing this ratio as $r\in\B$, we cannot have $r=1$ because $x\ne y$.  Hence
\[
        \alpha=\frac{rx-y}{r-1}\in\B,
\]
contradicting $\alpha\notin\B$.  Thus the pole vector is not projectively
$\B$-valued.  In particular the displayed affine parametrization has no
$\B$-valued direction vector after rescaling.
\end{proof}

\begin{corollary}[extension-pole quotient-remainder floor]
\label{cor:extension-pole-quotient-remainder-floor}
Assume $\B\subsetneq\F$, $D\subseteq\B^\times$ is a multiplicative coset of order
$n\ge2$, put $b:=|\B|$, and let $C=\RS[\F,D,k]$.
Fix a quotient-remainder profile as in
\cref{lem:heaviest-prefix-locator-floor} with $K=k+1$ and agreement $A_0>k$.
Let
\[
        L:=H_{c,m,s}^{k+1}
\]
be the certified heaviest prefix-fiber size, or replace $L$ by either fallback
\[
        \left\lceil \frac{M_{c,m,s}}{I_{c,m,s}^{k+1}}\right\rceil,
        \qquad
        \left\lceil\frac{M_{c,m,s}}{|\B|^{w_c(s,A_0-k-1)}}\right\rceil .
\]
Then, for every
\[
        \delta\in[1-A_0/n,1-k/n),
\]
there is a genuinely extension-valued simple-pole line with at least
\[
        \mathcal N^{\rm ext}_{q,b,k}(L)
        :=
        \left\lceil\frac{L(q-b)}{q-b+kL}\right\rceil
\]
CA-bad finite slopes.  For $\delta>0$, these slopes are MCA-bad.  Equivalently,
\[
        \eca(C,\delta)
        \ge
        \frac{L(q-b)}{q(q-b+kL)},
        \qquad
        \emca(C,\delta)
        \ge
        \frac{L(q-b)}{q(q-b+kL)}
        \quad(\delta>0).
\]
If $L>(q-b)/k$, this gives the clean trigger
\[
        \eca(C,\delta)>
        \frac1{2k}\left(1-\frac bq\right),
        \qquad
        \emca(C,\delta)>
        \frac1{2k}\left(1-\frac bq\right)
        \quad(\delta>0).
\]
\end{corollary}

\begin{proof}
\Cref{lem:heaviest-prefix-locator-floor} gives a $\B$-valued received word with
at least $L$ codewords of $C^+$ in the closed ball of radius $1-A_0/n$.  Apply
\cref{cor:extension-pole-deep-list-floor}.  The trigger inequality is
\[
        \frac{L(q-b)}{q(q-b+kL)}>
        \frac{q-b}{2kq},
\]
which is equivalent to $kL>q-b$.
\end{proof}

\begin{definition}[extension-pole certificate]
\label{def:extension-pole-certificate}
An extension-pole certificate for a list $P_1,\ldots,P_L\in\F[X]_{<k+1}$ and a
received word $U$ consists of:
\begin{enumerate}[label=\textup{(\roman*)},leftmargin=2.2em]
\item a declared proper subfield $\B\subsetneq\F$ containing $D$ and a pole
$\alpha\in\F\setminus\B$;
\item the agreement supports showing that each $P_i$ lies in the declared ball
around $U$;
\item the bucket table for the values $P_i(\alpha)$, with duplicate values
coalesced;
\item the simple-pole line $(f_\alpha,g_\alpha)$ and the proof that
$g_\alpha$ is not projectively $\B$-valued;
\item the final bad-slope numerator, namely the number of distinct buckets in
\textup{(iii)}.
\end{enumerate}
The verifier checks $\alpha\notin\B$, the support agreements, the value buckets,
and the simple-pole far condition.  For quotient-remainder floors, the list in
\textup{(ii)} may itself be supplied by a prefix-fiber certificate from
\cref{def:prefix-fiber-certificate}.
\end{definition}

\begin{remark}[what the extension-pole milestone changes]
\label{rem:milestone3-extension-status}
The subfield-confinement theorem says that $\B$-rational lines are inert over a
proper extension.  \Cref{cor:extension-pole-deep-list-floor} gives the matching
existence theorem on the failure side: any large list can be converted using
poles restricted to $\F\setminus\B$, so the resulting bad line is genuinely
extension-valued.  Thus the extension-line bucket in a prize scanner is no
longer only an abstract warning; it has an explicit simple-pole source and a
stable numerator with $q-b$ replacing $q-n$.  What remains open is not existence
of extension-valued witnesses, but their sharp classification inside the
aperiodic/residue-line ledgers and, for a fully explicit counterexample, printing
a concrete pole $\alpha$ with its value-bucket table.
\end{remark}

\begin{theorem}[{\cite[Thm.~1.9]{BCHKS25}}, as stated in {\cite[Thm.~5.2]{ABF26}}]\label{thm:B}
Let $C=\RS[\F,D,k]$, let $\rho=k/n$, and let $\delta\in(0,1-\rho)$ satisfy $\delta+2/n\le 1-\rho-1/n$. If
\[
\eca\Bigl(C,\ \delta+\tfrac2n,\ 1-\rho-\tfrac1n\Bigr)\;<\;\frac1{2n},
\qquad\text{then}\qquad
\Lst(C,\delta)\;<\;|\F| .
\]
\end{theorem}

\begin{remark}[on imports and radius conventions]\label{rem:import}
\Cref{thm:A} is the only Crites--Stewart-type consequence used by the main cap, and it is now proved above by the simple-pole deep-point argument.  The important admissibility condition is the integer-radius condition
\[
        \lfloor \delta n\rfloor \le n-k-1 .
\]
Consequently, the no-loss CA lower bound proved below should be read only in this admissible range. The Proximity-Prize MCA cap is unaffected: we prove the lower bound at the endpoint
\[
        \delta_N:=1-\rho-\frac2N
\]
and then use support-wise MCA monotonicity to extend the MCA failure to all larger sub-capacity radii.

\Cref{thm:B} remains a separate slacked fallback imported through the BCHKS/ABF route. It is not needed for the main field-size-universal MCA cap.
\end{remark}

\section{Locator fibers and interleaving transfer}\label{sec:fiber}

Fix $N\mid n$, set $a:=n/N$, and let $Q:=D^{a}:=\{x^a:x\in D\}\subseteq\B^\times$, a multiplicative coset of order $N$. The $a$-th power map $D\to Q$ is surjective and its fibers $S_b:=\{x\in D:x^a=b\}$ have exactly $a$ elements each. Since $a\mid n\mid|\B|-1$, the characteristic of $\F$ does not divide $a$, so $X^a-b$ is separable for $b\ne0$ and its full root set inside $D$ is $S_b$ whenever $b\in Q$. For a set $A$, $e_j(A)$ denotes the $j$-th elementary symmetric function of its elements.

\begin{lemma}[locator fibers are lists]\label{lem:fiber}
Let $\B\subseteq\F$, let $D\subseteq\B^\times$ be a multiplicative coset of order $n$, and let $k$ satisfy $a\mid k$ and $\rho=k/n$. Write $\ell_1:=\rho N+1$ and $\ell_2:=\rho N+2$.
\begin{enumerate}[label=\textup{(\roman*)}]
\item If $\ell_1\le N$, there exists $z\in\B$ such that, with $u_z:=\bigl(x^{k+a}+z\,x^{k}\bigr)_{x\in D}\in\F^n$,
\[
\bigl|\Lst\bigl(\RS[\F,D,k],\ 1-\rho-\tfrac1N,\ u_z\bigr)\bigr|\;\ge\;\binom{N}{\ell_1}\Big/|\B| .
\]
\item If $\ell_2\le N$, there exists $z\in\B$ such that, with $u_z:=\bigl(x^{k+2a}+z\,x^{k+a}\bigr)_{x\in D}\in\F^n$,
\[
\bigl|\Lst\bigl(\RS[\F,D,k+1],\ 1-\rho-\tfrac2N,\ u_z\bigr)\bigr|\;\ge\;\binom{N}{\ell_2}\Big/|\B| .
\]
\end{enumerate}
In both cases the words $u_z$ are $\B$-valued and all listed codewords lie in $\RS[\B,D,k]$ \textup(resp.\ $\RS[\B,D,k+1]$\textup).
\end{lemma}

\begin{proof}
We prove (ii); part (i) is the identical computation one degree rung lower and is \cite[Main Thm.\,(d)]{Cho26a} sharpened by pigeonholing over $\B$ in place of $\F$.

Let $A\subseteq Q$ with $|A|=\ell_2$, and put
\[
L_A(X)\;:=\;\prod_{b\in A}\bigl(X^a-b\bigr)\;=\;\sum_{j=0}^{\ell_2}(-1)^j e_j(A)\,X^{a(\ell_2-j)}\;=\;X^{k+2a}-e_1(A)\,X^{k+a}+R_A(X),
\]
where $a\ell_2=k+2a$, $a(\ell_2-1)=k+a$, and $R_A$ collects the terms with $j\ge2$, so $\deg R_A\le a(\ell_2-2)=k$. The polynomial $L_A$ is squarefree with root set $S_A:=\bigcup_{b\in A}S_b\subseteq D$ of size $a\ell_2=k+2a$.

Set $z_A:=-e_1(A)\in\B$ and $c_A:=\bigl(-R_A(x)\bigr)_{x\in D}$. Since $\deg R_A\le k$, we have $c_A\in\RS[\B,D,k+1]\subseteq\RS[\F,D,k+1]$. For every $x\in S_A$,
\[
0=L_A(x)=x^{k+2a}+z_A\,x^{k+a}+R_A(x),\qquad\text{i.e.}\qquad u_{z_A}(x)=c_A(x),
\]
so $u_{z_A}$ and $c_A$ agree on at least $k+2a$ points of $D$ and $\dist(u_{z_A},c_A)\le1-\rho-2/N$.

The map $A\mapsto c_A$ is injective on the $\ell_2$-subsets of $Q$ sharing a common slope value: if $z_A=z_{A'}$ and $c_A=c_{A'}$, then $R_A$ and $R_{A'}$ are polynomials of degree $\le k<n$ agreeing on all $n$ points of $D$, hence equal as polynomials; then $L_A=L_{A'}$, so $S_A=S_{A'}$, and $A=\{x^a:x\in S_A\}=A'$.

Finally, $z_A=-e_1(A)$ takes at most $|\B|$ values as $A$ ranges over the $\binom{N}{\ell_2}$ subsets; some $z\in\B$ is attained by at least $\binom{N}{\ell_2}/|\B|$ of them, and the corresponding codewords $c_A$ are pairwise distinct elements of $\Lst(\RS[\F,D,k+1],\,1-\rho-2/N,\,u_z)$.
\end{proof}

\begin{remark}[why slack two]\label{rem:slacktwo}
Part (i) applied to $C^+=\RS[\F,D,k+1]$ directly would require $a\mid k+1$; in every challenge instantiation $a$ is a power of two and $k+1$ is odd, so this fails. At slack two the listed codewords have degree $\le k$, i.e.\ lie in $C^+$, while the divisibility demand stays on $k$. This is exactly the interface needed by \cref{thm:A}, whose list conclusion concerns $C^+$ while its hypothesis concerns $C$.
\end{remark}

\begin{remark}[the subfield pigeonhole elsewhere]\label{rem:subpigeon}
The same one-line strengthening --- pigeonholing locator data over the field of definition $\B$ rather than over $\F$ --- also upgrades the coefficient-pigeonhole bound of \cite{Cho26b} to the generated field: the prefix map $\Phi_\sigma$ there has image in $\B^\sigma$ whenever $D\subseteq\B$, since the elementary symmetric prefixes of subsets of $D$ are $\B$-valued. This is the form consumed by the ``necessary'' half of the list-profile conjecture of \cite{Cho26b}, which states the entropy floor at the generated field $\tau^{*}(\rho,q_D)$; \cref{lem:fiber} supplies the missing subfield step.
\end{remark}

\begin{remark}[lists, not lines]\label{rem:lines}
The words $u_z$ of \cref{lem:fiber} lie on the $\B$-rational line $f+zg$ with $f=(x^{k+ta})_{x\in D}$, $g=(x^{k+(t-1)a})_{x\in D}$, but the lemma is purely a statement about lists at a received word and no line structure is used in \cref{sec:main}. This matters: by \cref{cor:Fvalued} below, over a proper extension $\F\supsetneq\B$ the lines that end up certifying the CA failure cannot be these (or any) $\B$-rational lines.
\end{remark}

\subsection{Quotient-profile floors with remainder supports}
\label{sec:quotient-remainder}

The fiber lemma above uses supports that are unions of complete quotient fibers.
For threshold work one also needs agreements that are not multiples of the fiber
size.  The following locator-prefix pigeonhole gives a finite lower ledger for
such remainder supports.  It is a lower-floor theorem only; the matching upper
ledger over all divisor scales remains part of the prize-completion program.

Let $c\mid n$ and put $Q_c:=D^c$, so the map
\[
        \pi_c:D\to Q_c,\qquad x\mapsto x^c
\]
has fibers of size $c$ and $|Q_c|=N=n/c$.  For $0\le s<c$ and an integer
$\sigma\ge0$, define the prefix weight
\[
        w_c(s,\sigma)
        :=\#\{h:1\le h\le\sigma,
              \ h\bmod c\in\{0,1,\ldots,s\}\},
\]
where residues are taken in $\{0,1,\ldots,c-1\}$.  Thus
$w_c(0,\sigma)=\lfloor \sigma/c\rfloor$ for full-fiber supports.

\begin{lemma}[quotient-remainder locator-prefix lower floor]
\label{lem:quotient-remainder-prefix}
Let $\B\subseteq\F$, let $D\subseteq\B^\times$ be a multiplicative coset of
order $n$, and let $K<n$.  Fix $c\mid n$, write $N=n/c$, and let
\[
        A_0=mc+s,
        \qquad 0\le s<c,
        \qquad 0\le m\le N,
        \qquad A_0\ge K .
\]
If $s>0$, assume $mc+s\le n$.  Put $\sigma=A_0-K$ and
\[
        M_{c,m,s}:=\binom Nm\binom{n-mc}{s}.
\]
Then there is a $\B$-valued received word $U:D\to\B$ such that
\[
        \bigl|\Lst(\RS[\F,D,K],\,1-A_0/n,\,U)\bigr|
        \ge
        \left\lceil \frac{M_{c,m,s}}{|\B|^{w_c(s,\sigma)}}\right\rceil .
\]
The listed codewords may be chosen in $\RS[\B,D,K]$.
\end{lemma}

\begin{proof}
For $M\subseteq Q_c$ with $|M|=m$, write
\[
        P_M(X)=\prod_{b\in M}(X^c-b).
\]
For $R\subseteq D\setminus \pi_c^{-1}(M)$ with $|R|=s$, put
\[
        E_R(X)=\prod_{r\in R}(X-r),
        \qquad
        L_{M,R}(X)=P_M(X)E_R(X).
\]
Then $L_{M,R}$ is a squarefree locator of degree $A_0$ whose roots are exactly
\(\pi_c^{-1}(M)\cup R\).  Write
\[
        L_{M,R}(X)=X^{A_0}+\sum_{h=1}^{A_0}\lambda_h(M,R)X^{A_0-h}.
\]
The coefficients lie in $\B$.  Moreover, because $P_M$ has monomials only in
degrees congruent to $0$ modulo $c$ and $E_R$ has degree $s$, the coefficient
$\lambda_h(M,R)$ is zero unless $h\bmod c\in\{0,1,\ldots,s\}$.
Consequently the high-degree prefix
\[
        \bigl(\lambda_h(M,R):1\le h\le\sigma,
              \ h\bmod c\in\{0,1,\ldots,s\}\bigr)
\]
takes at most $|\B|^{w_c(s,\sigma)}$ values.

Pigeonhole over these prefix values.  There is a fixed prefix vector $\xi$ and a
subfamily $\mathcal T_\xi$ of at least
$M_{c,m,s}/|\B|^{w_c(s,\sigma)}$ pairs $(M,R)$ with this prefix.  Let
\[
        U_\xi(X)=X^{A_0}+\sum_{1\le h\le\sigma}\lambda_h(\xi)X^{A_0-h}
\]
be the corresponding high-degree part, restricted to $D$.  For every
$(M,R)\in\mathcal T_\xi$, the difference
\[
        R_{M,R}(X):=L_{M,R}(X)-U_\xi(X)
\]
has degree $<K$.  Hence the codeword $-R_{M,R}$ agrees with $U_\xi$ on all
roots of $L_{M,R}$, a set of size $A_0$.

It remains only to check that the listed codewords are distinct.  If two pairs
$(M,R)$ and $(M',R')$ in the same prefix class gave the same codeword on $D$,
then the two remainders of degree $<K<n$ would be equal as polynomials.  Since
the high prefixes are also equal, $L_{M,R}=L_{M',R'}$.  Their root sets in $D$
are therefore equal.  A full fiber has size $c$ while $|R|=s<c$, so the full
fibers contained in this root set are exactly $M$, and then the residual set is
forced as well.  Thus $(M,R)=(M',R')$.
\end{proof}

\begin{definition}[prefix fibers and prefix-image certificates]
\label{def:prefix-fiber-certificate}
In the notation of \cref{lem:quotient-remainder-prefix}, let
$\Phi_{c,m,s}^{K}$ be the map from pairs $(M,R)$ to the high-degree prefix
\[
        \bigl(\lambda_h(M,R):1\le h\le A_0-K,
              \ h\bmod c\in\{0,1,\ldots,s\}\bigr)
        \in \B^{w_c(s,A_0-K)} .
\]
If $A_0=K$ this is the unique empty prefix.  For a prefix value $\xi$, put
\[
        \mu_{c,m,s}^{K}(\xi)
        :=\bigl| (\Phi_{c,m,s}^{K})^{-1}(\xi)\bigr|.
\]
Define the image size and the heaviest prefix fiber by
\[
        I_{c,m,s}^{K}:=\bigl|\operatorname{im}\Phi_{c,m,s}^{K}\bigr|,
        \qquad
        H_{c,m,s}^{K}:=\max_{\xi}\mu_{c,m,s}^{K}(\xi).
\]
Thus
\[
        H_{c,m,s}^{K}
        \ge
        \left\lceil\frac{M_{c,m,s}}{I_{c,m,s}^{K}}\right\rceil
        \ge
        \left\lceil\frac{M_{c,m,s}}{|\B|^{w_c(s,A_0-K)}}\right\rceil .
\]
A prefix-fiber certificate is an orbit decomposition, hash table, or other
verifiable bucket table for the finite map $\Phi_{c,m,s}^{K}$, including the
image size and the maximum fiber size.  If only $I_{c,m,s}^{K}$ is certified,
the middle lower bound is still valid; if no image certificate is supplied, the
ambient-box lower bound is valid.
\end{definition}

\begin{lemma}[heaviest-prefix locator floor]
\label{lem:heaviest-prefix-locator-floor}
Under the hypotheses of \cref{lem:quotient-remainder-prefix}, there is a
$\B$-valued received word $U:D\to\B$ such that
\[
        \bigl|\Lst(\RS[\F,D,K],\,1-A_0/n,\,U)\bigr|
        \ge H_{c,m,s}^{K}.
\]
Consequently
\[
        \bigl|\Lst(\RS[\F,D,K],\,1-A_0/n,\,U)\bigr|
        \ge
        \left\lceil \frac{M_{c,m,s}}{I_{c,m,s}^{K}}\right\rceil
        \ge
        \left\lceil \frac{M_{c,m,s}}{|\B|^{w_c(s,A_0-K)}}\right\rceil
\]
when only the coarser certificates are available.  The listed codewords may be
chosen in $\RS[\B,D,K]$.
\end{lemma}

\begin{proof}
Choose a prefix value $\xi$ with
$\mu_{c,m,s}^{K}(\xi)=H_{c,m,s}^{K}$ and form the high-degree word
\[
        U_\xi(X)=X^{A_0}+\sum_{1\le h\le A_0-K}
        \lambda_h(\xi)X^{A_0-h},
\]
where $\lambda_h(\xi)=0$ for the high coefficients outside the allowed
residue classes.  For every pair $(M,R)$ in this heaviest prefix fiber, the difference
$L_{M,R}(X)-U_\xi(X)$ has degree $<K$, so the negative of this remainder is a
codeword of $\RS[\B,D,K]$ agreeing with $U_\xi$ on the $A_0$ roots of
$L_{M,R}$.  If two pairs in the same prefix fiber gave the same codeword on
$D$, then their degree-$<K<n$ remainders would be equal as polynomials; together
with the common prefix this gives $L_{M,R}=L_{M',R'}$.  Their root sets in $D$
are equal.  A full quotient fiber has size $c$ while the residual set has size
$s<c$, so the full fibers contained in the root set determine $M$, and then the
residual set determines $R$.  Thus the codewords are pairwise distinct.  The two
coarser displayed bounds follow from \cref{def:prefix-fiber-certificate}.
\end{proof}

For $A\in\{k+1,
\ldots,n\}$ and $c\mid n$, write
\[
        A=m_c(A)c+s_c(A),
        \qquad 0\le s_c(A)<c,
\]
and set
\[
\begin{aligned}
        M_c(A)&:=
        \binom{n/c}{m_c(A)}
        \binom{n-cm_c(A)}{s_c(A)},\\
        L_{c,{\rm heavy}}^+(A)&:=
        H_{c,m_c(A),s_c(A)}^{k+1},\\
        L_{c,{\rm image}}^+(A)&:=
        \left\lceil\frac{M_c(A)}{I_{c,m_c(A),s_c(A)}^{k+1}}\right\rceil,\\
        L_{c,{\rm box}}^+(A)&:=
        \left\lceil
        \frac{M_c(A)}{|\B|^{w_c(s_c(A),A-k-1)}}
        \right\rceil .
\end{aligned}
\]
The three levels satisfy
\[
        L_{c,{\rm heavy}}^+(A)
        \ge L_{c,{\rm image}}^+(A)
        \ge L_{c,{\rm box}}^+(A).
\]
Here ``heavy'' uses a full prefix-fiber certificate, ``image'' uses only the
prefix-image size, and ``box'' is the certificate-free bound already implicit in
\cref{lem:quotient-remainder-prefix}.

\begin{theorem}[quotient-remainder deep-point lower ledger]
\label{thm:quotient-remainder-deep-floor}
Let $\B\subseteq\F$, let $q=|\F|>n$, let
$D\subseteq\B^\times$ be a multiplicative coset of order $n$, and let
$C=\RS[\F,D,k]$.  Fix an agreement value $A\in\{k+1,\ldots,n\}$ and a
divisor $c\mid n$.  For each
\[
        \star\in\{{\rm heavy},{\rm image},{\rm box}\}
\]
and every
\[
        \delta\in[1-A/n,\,1-k/n)
\]
one has
\[
        \eca(C,\delta)
        \ge \mathcal E_{q,k}\bigl(L_{c,\star}^+(A)\bigr).
\]
If $\delta>0$, then also
\[
        \emca(C,\delta)
        \ge \mathcal E_{q,k}\bigl(L_{c,\star}^+(A)\bigr).
\]
Here $\mathcal E_{q,k}$ is as in \cref{cor:deep-list-trigger-ceiling}.  More
precisely, some received line has at least
\[
        \mathcal N_{q,k}\bigl(L_{c,\star}^+(A)\bigr)
\]
CA-bad finite slopes at that radius, and for $\delta>0$ these slopes are
MCA-bad.

Consequently, for a closed agreement threshold $a\ge k+1$, define
\[
        B_{Q,\star}^-(a)
        :=\max_{\substack{c\mid n\\ A\in\{a,a+1,\ldots,n\}}}
          \mathcal N_{q,k}\bigl(L_{c,\star}^+(A)\bigr).
\]
Then $B_{Q,\star}^-(a)$ is a lower numerator for $q\,\eca(C,\delta)$ for
every $\delta\in[1-a/n,1-k/n)$, and for $q\,\emca(C,\delta)$ on the
same interval with $\delta>0$.  The strongest
certified ledger uses $\star={\rm heavy}$; the certificate-free ledger uses
$\star={\rm box}$.
\end{theorem}

\begin{proof}
Apply \cref{lem:heaviest-prefix-locator-floor} with $K=k+1$ and $A_0=A$.  This
produces a $C^+=\RS[\F,D,k+1]$ list of size at least
$L_{c,{\rm heavy}}^+(A)$ at radius $1-A/n$, with the image and ambient-box
variants giving the corresponding smaller certified list sizes.  Then apply
\cref{prop:quantitative-deep-list-floor}.  If $A\ge a$ and
$\delta\ge1-a/n$, then $\delta\ge1-A/n$, so the same converted floor applies.
The maximum over $(c,A)$ is the correct lower-ledger operation: the definitions
of $\eca$ and $\emca$ maximize over received lines, so lower witnesses from
different divisor scales or exact agreement values need not share one line and
must not be summed.
\end{proof}

\begin{corollary}[quotient-remainder trigger]
\label{cor:quotient-remainder-trigger}
In the setting of \cref{thm:quotient-remainder-deep-floor}, if
$L_{c,\star}^+(A)>(q-n)/k$ for any
$\star\in\{{\rm heavy},{\rm image},{\rm box}\}$, then, for every
$\delta\in[1-A/n,1-k/n)$,
\[
        \eca(C,\delta)>
        \frac1{2k}\left(1-\frac nq\right),
\]
and for every such $\delta>0$,
\[
        \emca(C,\delta)>
        \frac1{2k}\left(1-\frac nq\right).
\]
\end{corollary}

\begin{proof}
This is \cref{cor:deep-list-trigger-ceiling} applied inside
\cref{thm:quotient-remainder-deep-floor}.
\end{proof}

\begin{corollary}[quantitative first-grid floor]
\label{cor:quantitative-first-grid-floor}
Let $C=\RS[\F,D,k]$, let $q=|\F|>n$, and assume only that $D\subseteq\F$ has
$n$ distinct points.  Then for every
$\delta\in[1-(k+1)/n,1-k/n)$,
\[
        \eca(C,\delta)
        \ge
        \mathcal E_{q,k}\left(\binom n{k+1}\right),
\]
and for every such $\delta>0$,
\[
        \emca(C,\delta)
        \ge
        \mathcal E_{q,k}\left(\binom n{k+1}\right).
\]
In particular, if $\binom n{k+1}>(q-n)/k$, then the printed
$(2k)^{-1}(1-n/q)$ lower bound holds on this entire first grid interval without
any multiplicative smoothness assumption.
\end{corollary}

\begin{proof}
Use the ordinary locator identity for all $(k+1)$-subsets of $D$; this is the
$c=1$ argument behind \cref{lem:heaviest-prefix-locator-floor}, but it does not
use multiplicative smoothness.  The prefix is empty, so the list size is at
least
$\binom n{k+1}$.  Apply \cref{prop:quantitative-deep-list-floor} and then
\cref{cor:deep-list-trigger-ceiling}.
\end{proof}

\begin{definition}[quotient lower-ledger certificate]
\label{def:quotient-lower-ledger-certificate}
A quotient-remainder lower-ledger certificate for a row and a threshold $a$
contains, for each divisor scale $c$ and each exact agreement $A\ge a$ that is
used in the maximum:
\begin{enumerate}[label=\textup{(\roman*)},leftmargin=2.2em]
\item the decomposition $A=m_c(A)c+s_c(A)$;
\item the support count $M_c(A)$ and prefix weight $w_c(s_c(A),A-k-1)$;
\item either a full prefix-fiber bucket table certifying
$H_{c,m_c(A),s_c(A)}^{k+1}$, or an image-size certificate for
$I_{c,m_c(A),s_c(A)}^{k+1}$, or the declaration that the ambient-box denominator
is being used;
\item the resulting list floor $L_{c,\star}^+(A)$;
\item the converted slope numerator
$\mathcal N_{q,k}(L_{c,\star}^+(A))$ and error floor
$\mathcal E_{q,k}(L_{c,\star}^+(A))$.
\end{enumerate}
The verifier checks the locator-prefix formula, distinctness of codewords inside
a prefix fiber, the integer arithmetic, and the final maximum over divisor
scales and exact agreement values.
\end{definition}

\begin{remark}[status of T1 after the v10 lower-ledger upgrade]
\label{rem:T1-lower-ledger-status-v10}
The quotient-profile lower side is now scanner-ready at every agreement value
$A=mc+s$: the generated-field contribution is the actual heaviest prefix fiber
$H_{c,m,s}^{k+1}$ when certified, the exact prefix-image denominator
$I_{c,m,s}^{k+1}$ when only image data are certified, and the ambient box
$|\B|^{w_c(s,A-k-1)}$ when no certificate is supplied.  The challenge-field size
$q$ enters only through the independent conversion functions
$\mathcal E_{q,k}$ and $\mathcal N_{q,k}$.  Across divisor scales the lower-side
operation is a maximum, not a union count, because the witness line may change.
The remaining T1 upper-side problem is sharper and unchanged in logical status:
replace the safe support ledger and the exact declared-branch image ledger by an
exhaustive quotient-periodic classification, with duplicate slopes coalesced
across all divisor scales.
\end{remark}

\begin{corollary}[augmented slack-one quotient floor]
\label{cor:augmented-slack-one}
Let $\B\subseteq\F$ and let $D\subseteq\B^\times$ be a multiplicative coset
of order $n$.  Put $C=\RS[\F,D,k]$ and $C^+=\RS[\F,D,k+1]$.  If
$c\mid\gcd(n,k)$, then there is a $\B$-valued received word $U_c$ such that
\[
        \bigl|\Lst(C^+,\,1-(k+c)/n,\,U_c)\bigr|
        \ge
        \binom{n/c}{k/c+1} .
\]
One may take $U_c(x)=x^{k+c}$.  In particular, for $c=1$ this gives the
universal first-grid floor
\[
        \bigl|\Lst(C^+,\,1-(k+1)/n,\,x^{k+1})\bigr|
        \ge
        \binom n{k+1},
\]
with no quotient denominator.  The $c=1$ statement in fact holds for any set
$D\subseteq\B$ of $n$ distinct points, not only for multiplicative cosets.
\end{corollary}

\begin{proof}
For $c>1$, apply Lemma~\ref{lem:quotient-remainder-prefix} with $K=k+1$,
$s=0$, and $m=k/c+1$.  Then $A_0=k+c$ and $\sigma=A_0-K=c-1$, so
$w_c(0,c-1)=0$.  The high-degree prefix is only the monomial $X^{k+c}$, giving
the displayed word.  For $c=1$, the same argument is just the ordinary locator
identity for all $(k+1)$-subsets of $D$: the high-degree prefix of a monic
locator of degree $k+1$ is $X^{k+1}$, and its remainder has degree at most $k$.
\end{proof}

\begin{corollary}[first-grid deep-point cap]
\label{cor:first-grid-cap}
Let $C=\RS[\F,D,k]$, $q=|\F|$, and suppose $q>n$.  Let
$c\mid\gcd(n,k)$.  If $c>1$, assume $D\subseteq\B^\times$ is a
multiplicative coset as above; for $c=1$, no smoothness is needed.  If
\[
        \binom{n/c}{k/c+1}\ge q/k+1,
\]
then
\[
        \eca\bigl(C,1-(k+c)/n\bigr)
        >\frac1{2k}\left(1-\frac nq\right),
        \qquad
        \emca\bigl(C,1-(k+c)/n\bigr)
        >\frac1{2k}\left(1-\frac nq\right).
\]
For $c=1$ this is the first closed grid point below capacity.
\end{corollary}

\begin{proof}
Corollary~\ref{cor:augmented-slack-one} gives a $C^+$-list of size
$L\ge q/k+1$ at agreement $k+c$.  The proof of the deep-point conversion
Theorem~\ref{thm:A} applies whenever the agreement is $>k$, so it applies here
for every $c\ge1$.  If the displayed CA lower bound failed, then
\[
        \eca(C,1-(k+c)/n)\le \frac{q-n}{2kq}.
\]
Theorem~\ref{thm:A} with $\eta=1/2$ would give
\[
        \Lst(C^+,1-(k+c)/n)
        \le
        \left\lceil \frac{q-n}{k}\right\rceil
        < q/k+1,
\]
contradicting the list lower bound.  Every no-loss CA-bad slope is also
support-wise MCA-bad on the same explaining support, giving the MCA inequality.
\end{proof}

\begin{corollary}[first-grid cap for large challenge-envelope rows]
\label{cor:first-grid-grand}
Let $\rho\in\{\frac12,\frac14,\frac18,\frac1{16}\}$, let
$C=\RS[\F,D,k]$ with $|D|=n$, $k=\rho n$, $q=|\F|<2^{256}$, and
$q>n$.  Assume
\[
\begin{array}{c|cccc}
\rho & \frac12 & \frac14 & \frac18 & \frac1{16}\\ \hline
k\ge & 127 & 78 & 58 & 47
\end{array}
\]
for the corresponding rate.  Then the first closed grid point below capacity is
already CA- and MCA-unsafe:
\[
        \eca\bigl(C,1-\rho-1/n\bigr)
        >\frac1{2k}\left(1-\frac nq\right),
        \qquad
        \emca\bigl(C,1-\rho-1/n\bigr)
        >\frac1{2k}\left(1-\frac nq\right).
\]
If also $k\le2^{40}$, then this lower bound is $>2^{-86}$ and hence is far
above the prize target $2^{-128}$.
\end{corollary}

\begin{proof}
By Corollary~\ref{cor:first-grid-cap} with $c=1$, it is enough to check
\[
        \binom n{k+1}\ge q/k+1.
\]
Since $q<2^{256}$, the stronger sufficient condition is
$k(\binom n{k+1}-1)\ge2^{256}$.  The four printed lower bounds on $k$ are exactly
small finite checks of this inequality at the corresponding rate; the left side
is increasing thereafter.  For the error size, $q\ge n+1$, $n=k/\rho\le16k$,
and $k\le2^{40}$ give
\[
        \frac1{2k}\left(1-\frac nq\right)
        \ge \frac1{2k(n+1)}
        >2^{-86}.
\]
\end{proof}


\subsection{A conservative quotient-support upper ledger}
\label{sec:quotient-support-upper}

The remainder lower floor above supplies the unsafe side of the quotient-profile
ledger.  For a prize threshold one also needs a safe-side accounting rule: if a
later argument declares a bad parameter to be quotient-periodic, the declaration
must come with a divisor-lattice numerator that is counted with overlaps under
control.  The next statements give such a conservative upper ledger.  They do
not prove that every bad line parameter is quotient-periodic; they only say that
once a permitted support family has been specified, every ledger term witnessed
inside that family is bounded by the size of the family, with a factor $d$ for
finite-parameter degree-$d$ power curves.

Let $\mathcal S$ be a family of subsets of $D$, all of size at least $a$.  For a
received line $f+zg$, say that a finite slope is
\emph{$\mathcal S$-witnessed support-wise bad} if it has a support-wise
noncontained explanation on some $S\in\mathcal S$.  Define the analogous
$\mathcal S$-witnessed no-loss CA slopes by requiring that $f+zg$ be explained
on some $S\in\mathcal S$ while $(f,g)$ is not globally explained on any support
of size at least $a$.

\begin{lemma}[one support pays for at most one line parameter]
\label{lem:one-support-one-line}
Let $C=\RS[\F,D,k]$ and let $a\ge k$.  Fix a support
$S\subseteq D$ with $|S|\ge a$.  For any received line $f+zg$, at most one finite
slope $z\in\F$ can be support-wise noncontained on the fixed support $S$.
Likewise, at most one finite slope can be no-loss CA-bad while being explained on
this fixed support $S$.  Consequently, for every support family $\mathcal S$,
\[
        \#\{\mathcal S\text{-witnessed support-wise MCA-bad finite slopes}\}
        \le |\mathcal S|,
\]
and the same bound holds for no-loss CA slopes witnessed by $\mathcal S$.
\end{lemma}

\begin{proof}
Suppose two distinct finite slopes $z_1\ne z_2$ are both explained on the same
support $S$, by codewords $p_1,p_2\in C$:
\[
        f+z_i g=p_i\qquad\text{on }S,\quad i=1,2.
\]
Then on $S$,
\[
        g=\frac{p_1-p_2}{z_1-z_2},\qquad
        f=p_1-z_1g.
\]
Both right-hand sides are codewords of $C$.  Thus $(f,g)$ is simultaneously
explained on $S$.  This contradicts support-wise noncontainment on $S$.  It also
contradicts no-loss CA badness at agreement $a$, because $|S|\ge a$ gives a
global simultaneous explanation support of the forbidden size.  Hence each fixed
support contributes at most one bad finite line parameter, and summing over
supports gives the displayed bound.
\end{proof}

For a finite-parameter power curve
\[
        W_\gamma=f_0+\gamma f_1+\cdots+\gamma^d f_d,
        \qquad \gamma\in\F,
\]
say that a parameter is support-wise curve-bad on $S$ if $W_\gamma$ is explained
by a codeword on $S$, but the coefficient tuple $(f_0,\ldots,f_d)$ is not
simultaneously explained by $C^{\equiv(d+1)}$ on $S$.  The no-loss curve-CA
version replaces this support-wise noncontainment by the corresponding global
far condition at agreement $a$.

\begin{lemma}[one support pays for at most $d$ curve parameters]
\label{lem:one-support-d-curve}
Let $C=\RS[\F,D,k]$, let $a\ge k$, and fix a support $S\subseteq D$ with
$|S|\ge a$.  For a degree-$d$ power curve $W_\gamma$, at most $d$ finite
parameters $\gamma\in\F$ can be support-wise curve-bad on $S$.  The same bound
holds for no-loss curve-CA parameters explained on $S$.  Consequently any
support family $\mathcal S$ witnesses at most $d|\mathcal S|$ bad finite curve
parameters.
\end{lemma}

\begin{proof}
If $d+1$ distinct parameters $\gamma_0,\ldots,\gamma_d$ were explained on the
same support $S$, with codewords $P_j\in C$, then Lagrange interpolation in the
parameter variable gives codewords $Q_0,\ldots,Q_d\in C$ such that
\[
        \sum_{i=0}^d \gamma_j^i Q_i=P_j
        \qquad (0\le j\le d).
\]
On every point of $S$, the same interpolation applied to the values of
$W_{\gamma_j}$ gives $f_i=Q_i$ for every coefficient $i$.  Hence
$(f_0,\ldots,f_d)$ is simultaneously explained by $C^{\equiv(d+1)}$ on $S$.
This contradicts support-wise curve-noncontainment on $S$, and also contradicts
no-loss curve-CA badness at agreement $a$.  Thus every support contributes at
most $d$ parameters.
\end{proof}

\begin{lemma}[one support pays for at most one interleaved tuple]
\label{lem:one-support-one-interleaved}
Let $C=\RS[\F,D,k]$, let $a\ge k$, and let $\mu\ge1$.  Fix a $\mu$-row received
word $U=(U_1,\ldots,U_\mu)$.  For a fixed support $S\subseteq D$ with
$|S|\ge a$, there is at most one tuple $(P_1,\ldots,P_\mu)\in C^\mu$ agreeing
with $U$ on every row over $S$.  Hence the number of interleaved list tuples whose
common agreement support belongs to a support family $\mathcal S$ is at most
$|\mathcal S|$.
\end{lemma}

\begin{proof}
For each row $i$, two degree-$<k$ polynomials agreeing with $U_i$ on the fixed
support $S$ agree with each other on at least $k$ points, so they are equal.  Thus
each row has at most one codeword compatible with $S$, and therefore the whole
interleaved tuple is unique if it exists.
\end{proof}

We now instantiate the support-family ledger on quotient-remainder supports.  For
$c\mid n$ and an integer $A$ with $0\le A\le n$, write
\[
        A=m_c(A)c+s_c(A),\qquad 0\le s_c(A)<c,
\]
and define
\[
\begin{aligned}
        \mathcal S_c(A)
        :=\{&\pi_c^{-1}(M)\cup R:
          M\subseteq Q_c,
          |M|=m_c(A),\\
        & R\subseteq D\setminus\pi_c^{-1}(M),
          |R|=s_c(A)\} .
\end{aligned}
\]
For a divisor set $\mathcal C\subseteq\{c:c\mid n\}$ put
\[
        \mathcal S_{\mathcal C}^{\ge a}
        :=\bigcup_{c\in\mathcal C}\bigcup_{A=a}^n \mathcal S_c(A)
\]
and
\[
        U_{\mathcal C}^{\rm supp}(a)
        :=\left|\mathcal S_{\mathcal C}^{\ge a}\right|.
\]
When the exact union is inconvenient, use the safe sum
\[
\label{eq:quotient-support-safe-sum}
        U_{\mathcal C}^{\rm sum}(a)
        :=\sum_{c\in\mathcal C}\sum_{A=a}^n
          \binom{n/c}{\lfloor A/c\rfloor}
          \binom{n-c\lfloor A/c\rfloor}{A-c\lfloor A/c\rfloor},
\]
so that $U_{\mathcal C}^{\rm supp}(a)\le U_{\mathcal C}^{\rm sum}(a)$.

\begin{proposition}[divisor-union quotient support ledger]
\label{prop:divisor-union-support-ledger}
Let $C=\RS[\F,D,k]$ and let $a\ge k$.  Fix a divisor set $\mathcal C$ and let
$\mathcal S_{\mathcal C}^{\ge a}$ be the quotient-remainder support family above.
Then, for every received line, the number of finite support-wise MCA-bad slopes
whose witness support lies in $\mathcal S_{\mathcal C}^{\ge a}$ is at most
\[
        U_{\mathcal C}^{\rm supp}(a)
        \le U_{\mathcal C}^{\rm sum}(a).
\]
The same bound holds for no-loss CA slopes witnessed by this family.  For a
finite-parameter degree-$d$ power curve, the corresponding support-wise curve-MCA
and no-loss curve-CA numerator is at most
\[
        d\,U_{\mathcal C}^{\rm supp}(a)
        \le d\,U_{\mathcal C}^{\rm sum}(a).
\]
For every interleaving arity $\mu\ge1$, the number of interleaved list tuples whose
common agreement support belongs to $\mathcal S_{\mathcal C}^{\ge a}$ is at most
$U_{\mathcal C}^{\rm supp}(a)$.
\end{proposition}

\begin{proof}
Apply Lemmas~\ref{lem:one-support-one-line}, \ref{lem:one-support-d-curve}, and
\ref{lem:one-support-one-interleaved} to the support family
$\mathcal S_{\mathcal C}^{\ge a}$.  The inequality
$U_{\mathcal C}^{\rm supp}(a)\le U_{\mathcal C}^{\rm sum}(a)$ is only the union
bound over the divisor and agreement-size indices.
\end{proof}

\begin{remark}[what this does and does not finish]
\label{rem:quotient-upper-status}
Proposition~\ref{prop:divisor-union-support-ledger} supplies a certified
safe-side quotient ledger at the support-family level, and it is already enough
for a staircase scanner to avoid double-counting divisor scales: one may compute
the exact union $U_{\mathcal C}^{\rm supp}$ or use the printed safe sum.  This is
not yet the sharp quotient upper theorem needed for a full prize solution.  Near
capacity the support count can be exponentially larger than the number of
distinct bad slopes, so the next refinement is a slope-collision quotient ledger:
classify when many quotient-remainder supports map to the same line or curve
parameter, and replace the support count by the corresponding distinct-parameter
union count.
\end{remark}


\begin{definition}[block support-profile polynomial]
\label{def:block-support-profile-polynomial}
Let $D$ be a multiplicative coset of order $n$, let
$\mathcal C$ be a nonempty finite set of divisors of $n$, and put
\[
        h:=\operatorname{lcm}\mathcal C,
        \qquad N_h:=n/h .
\]
For each $c\in\mathcal C$ let $H_c\le H_h$ denote the subgroup of the cyclic
kernel of order $h$ with $|H_c|=c$.  If $T\subseteq H_h$, define
\[
        \phi_c(T):=\#\{\text{cosets }uH_c\subseteq H_h:\ uH_c\subseteq T\},
\]
the number of complete $c$-subfibers contained in the block subset $T$.
The block support-profile polynomial of $\mathcal C$ is
\[
        G_{\mathcal C,h}(Y,(Z_c)_{c\in\mathcal C})
        :=\sum_{T\subseteq H_h}
          Y^{|T|}\prod_{c\in\mathcal C}Z_c^{\phi_c(T)} .
\]
For a vector $\nu=(\nu_c)_{c\in\mathcal C}$ write
$Z^\nu:=\prod_{c\in\mathcal C}Z_c^{\nu_c}$, and define
$N_{\mathcal C,h}(A;\nu)$ by
\[
        G_{\mathcal C,h}(Y,Z)^{N_h}
        =\sum_{A=0}^{n}\sum_{\nu}
          N_{\mathcal C,h}(A;\nu)Y^A Z^\nu .
\]
\end{definition}

\begin{lemma}[block polynomial counts exact quotient profiles]
\label{lem:block-polynomial-counts-profiles}
For a support $S\subseteq D$ and $c\mid n$, let
\[
        F_c(S):=\#\{\text{$c$-fibers of }D\to D^c\text{ contained in }S\}.
\]
Then $N_{\mathcal C,h}(A;\nu)$ is exactly the number of supports
$S\subseteq D$ such that
\[
        |S|=A,
        \qquad F_c(S)=\nu_c\quad(c\in\mathcal C).
\]
\end{lemma}

\begin{proof}
The $h$-fibers partition $D$ into $N_h$ blocks, and every $c$-fiber with
$c\in\mathcal C$ is contained in a unique $h$-fiber because $c\mid h$.  After
choosing an identification of each $h$-fiber with the cyclic kernel $H_h$, the
number of selected points and the number of complete internal $c$-subfibers are
recorded by the monomial
\[
        Y^{|T|}\prod_{c\in\mathcal C}Z_c^{\phi_c(T)}
\]
for the block subset $T$.  The choices in distinct $h$-fibers are independent,
and the global statistics $|S|$ and $F_c(S)$ are the sums of the block statistics.
Multiplying the block polynomial over the $N_h$ blocks and extracting the
coefficient gives the asserted count.
\end{proof}

\begin{lemma}[membership by full-fiber profile]
\label{lem:quotient-support-membership-profile}
Fix $A$ and $c\mid n$, and write
\[
        A=m_c(A)c+s_c(A),
        \qquad 0\le s_c(A)<c .
\]
For a support $S\subseteq D$ of size $A$, one has
\[
        S\in\mathcal S_c(A)
        \qquad\Longleftrightarrow\qquad
        F_c(S)=m_c(A)=\lfloor A/c\rfloor .
\]
\end{lemma}

\begin{proof}
If $S\in\mathcal S_c(A)$, then by definition it is the union of
$m_c(A)$ complete $c$-fibers and $s_c(A)<c$ residual points outside those fibers;
there is no room for any further complete $c$-fiber.  Hence
$F_c(S)=m_c(A)$.  Conversely, if $F_c(S)=m_c(A)$, take the complete $c$-fibers
contained in $S$ as the set $M$ in the definition of $\mathcal S_c(A)$; the
remaining $A-m_c(A)c=s_c(A)$ selected points form the residual set.
\end{proof}

\begin{theorem}[exact divisor-lattice support-union formula]
\label{thm:exact-divisor-block-support-ledger}
Let $\mathcal C$ be a nonempty finite set of divisors of $n$, and define
$G_{\mathcal C,h}$ and $N_{\mathcal C,h}(A;\nu)$ as in
\cref{def:block-support-profile-polynomial}.  Then the exact support-family
union from \cref{prop:divisor-union-support-ledger} is
\[
\boxed{
        U_{\mathcal C}^{\rm supp}(a)
        =\sum_{A=a}^{n}\sum_{\nu}
          N_{\mathcal C,h}(A;\nu)\,
          \mathbf 1\!
          \left[\exists c\in\mathcal C:
                 \nu_c=\left\lfloor A/c\right\rfloor\right].
      }
\]
Consequently the quotient-support branch numerators in
\cref{prop:divisor-union-support-ledger} may be computed by this coefficient
formula instead of by the safe sum
$U_{\mathcal C}^{\rm sum}(a)$.
\end{theorem}

\begin{proof}
For a fixed exact size $A$, the union
$\bigcup_{c\in\mathcal C}\mathcal S_c(A)$ consists precisely of those supports
$S$ of size $A$ for which $S\in\mathcal S_c(A)$ for at least one divisor
$c\in\mathcal C$.  By \cref{lem:quotient-support-membership-profile}, this is
exactly the condition
\[
        F_c(S)=\lfloor A/c\rfloor
\]
for at least one $c\in\mathcal C$.  By
\cref{lem:block-polynomial-counts-profiles}, the number of supports with a fixed
profile vector $\nu=(F_c(S))_{c\in\mathcal C}$ is
$N_{\mathcal C,h}(A;\nu)$.  Summing over all exact sizes $A\ge a$ and over the
profile vectors satisfying the displayed predicate gives the formula.
\end{proof}

\begin{corollary}[single-divisor sanity check]
\label{cor:single-divisor-support-ledger}
For $\mathcal C=\{c\}$ the block polynomial is
\[
        G_{\{c\},c}(Y,Z)=(1+Y)^c+(Z-1)Y^c,
\]
and \cref{thm:exact-divisor-block-support-ledger} gives
\[
        \left|\mathcal S_c(A)\right|
        =\binom{n/c}{\lfloor A/c\rfloor}
          \binom{n-c\lfloor A/c\rfloor}{A-c\lfloor A/c\rfloor},
\]
the summand used in the safe ledger.
\end{corollary}

\begin{proof}
Inside one $c$-fiber, every proper subset contributes only $Y^{|T|}$, while the
full subset contributes $Y^cZ$.  Thus the block polynomial is
$(1+Y)^c-Y^c+Y^cZ$.  Raising to $n/c$ and selecting exactly
$m=\lfloor A/c\rfloor$ full blocks and $s=A-mc$ residual points in the remaining
blocks gives the displayed count.
\end{proof}

\begin{definition}[divisor-block support-ledger certificate]
\label{def:divisor-block-support-ledger-certificate}
A divisor-block support-ledger certificate for a divisor set $\mathcal C$ and a
threshold $a$ consists of:
\begin{enumerate}[label=\textup{(\roman*)},leftmargin=2.2em]
\item the lcm block size $h=\operatorname{lcm}\mathcal C$ and the block count
$N_h=n/h$;
\item the internal coset table of every subgroup $H_c\le H_h$ for
$c\in\mathcal C$;
\item the block polynomial $G_{\mathcal C,h}$, either as a full coefficient table
or as an orbit-compressed table under translations of $H_h$;
\item the coefficient extraction from $G_{\mathcal C,h}^{N_h}$ for every exact
size $A\ge a$;
\item the predicate table
$\mathbf 1[\exists c\in\mathcal C:\nu_c=\lfloor A/c\rfloor]$ and the final
integer $U_{\mathcal C}^{\rm supp}(a)$.
\end{enumerate}
The verifier checks the block enumeration, the subgroup containment relation
$c\mid h$, coefficient arithmetic, and the final predicate sum.  This is an exact
support-union certificate; it does not by itself coalesce distinct line or curve
parameters that arise from different supports.
\end{definition}

\begin{remark}[support milestone versus slope milestone]
\label{rem:support-vs-slope-milestone}
\Cref{thm:exact-divisor-block-support-ledger} closes the divisor-lattice overlap
problem at the support-family level: intersections between quotient descriptions
at incomparable divisor scales are now counted by one coefficient extraction in
the common $h$-block.  This is still weaker than the desired distinct-slope
quotient ledger.  To finish that stronger milestone one must apply the exact
support-image maps of
\cref{prop:distinct-parameter-line-ledger,prop:distinct-parameter-curve-ledger}
to the support set counted here and then prove or certify the resulting parameter
image size after duplicate coalescing.
\end{remark}

\subsection{Distinct-parameter quotient image ledger}
\label{sec:quotient-distinct-parameter}

The support-union ledger of \cref{sec:quotient-support-upper} is deliberately
safe but often very loose: many quotient-remainder supports can certify the same
line slope or the same curve parameter.  The next refinement is to count the
image of the actual parameter map.  This subsection gives an exact image ledger
for the quotient branch.  It is still a branch ledger, not an aperiodic upper
bound: it counts only witnesses whose agreement support lies in a declared
quotient-remainder support family.

We first record the support-to-parameter map in syndrome form.  Let
$C=\RS[\F,D,k]$, $|D|=n$, and put $r=n-k$.  For each $x\in D$ set
\[
        \lambda_x:=\left(\prod_{y\in D,\ y\ne x}(x-y)\right)^{-1}.
\]
For a word $Y:D\to\F$, define its syndrome vector
\[
        \Syn(Y)_m:=\sum_{x\in D}\lambda_x x^mY(x),
        \qquad 0\le m<r .
\]
Thus $\Syn(Y)=0$ if and only if $Y\in C$.  If
$T\subseteq D$ has size $j\le r$, write
\[
        L_T(X)=\prod_{x\in T}(X-x)=\ell_0+\ell_1X+\cdots+\ell_jX^j,
        \qquad
        \boldsymbol\ell_T=(\ell_0,\ldots,\ell_j)^t,
\]
and put $t=r-j$.  For $w\in\F^r$ let $H_{t,j}(w)$ be the $t\times(j+1)$ Hankel
window
\[
        H_{t,j}(w)_{a,b}=w_{a+b},
        \qquad 0\le a<t,\quad 0\le b\le j .
\]
When $S=D\setminus T$, the integer $t$ is exactly $|S|-k$.


\begin{lemma}[support-locator syndrome recurrence]
\label{lem:support-locator-syndrome-recurrence}
Let $Y:D\to\F$ and let $T\subseteq D$ with $|T|=j\le r$.  Put
$S=D\setminus T$ and $t=r-j$.  Then $Y$ is explained by $C$ on $S$ if and only if
\[
        H_{t,j}(\Syn(Y))\boldsymbol\ell_T=0 .
\]
\end{lemma}

\begin{proof}
Let $W_T\subseteq\F^r$ be the span of the parity-check columns indexed by $T$,
\[
        W_T=\operatorname{span}_\F\{\lambda_x(1,x,\ldots,x^{r-1}):x\in T\}.
\]
The word $Y$ is explained by $C$ on $S$ if and only if $Y-c$ is supported on $T$
for some $c\in C$, equivalently $\Syn(Y)\in W_T$.  If
$w\in W_T$, write $w_m=\sum_{x\in T}a_xx^m$.  Then, for $0\le h<t$,
\[
        \bigl(H_{t,j}(w)\boldsymbol\ell_T\bigr)_h
        =\sum_{b=0}^j\ell_b w_{h+b}
        =\sum_{x\in T}a_xx^hL_T(x)=0.
\]
Conversely, because $L_T$ is monic, the displayed recurrence determines all
coordinates $w_j,\ldots,w_{r-1}$ from the first $j$ coordinates, so its solution
space has dimension at most $j$.  The $j$ parity-check columns indexed by the
distinct points of $T$ have Vandermonde rank $j$, hence $W_T$ has dimension $j$
and is exactly the recurrence space.
\end{proof}

\begin{lemma}[exact affine and projective support-image map]
\label{lem:exact-line-image-map}
Let $f,g:D\to\F$, put $u=\Syn(f)$ and $v=\Syn(g)$, and fix
$T\subseteq D$ with $|T|=j\le r$ and $S=D\setminus T$.  Define
\[
        A_T:=H_{t,j}(u)\boldsymbol\ell_T,
        \qquad
        B_T:=H_{t,j}(v)\boldsymbol\ell_T .
\]
A finite slope $z\in\F$ is explained on the fixed support $S$ and is
support-wise noncontained there if and only if
\[
        A_T+zB_T=0,
        \qquad B_T\ne0 .
\]
Consequently the fixed support contributes either no finite bad slope or the
unique value
\[
        \theta_T=-A_{T,m}/B_{T,m}
\]
for any coordinate $m$ with $B_{T,m}\ne0$, provided all nonzero coordinates give
the same value.

For the projective line $\alpha f+\beta g$, the fixed support contributes a
support-wise noncontained point $[\alpha:\beta]\in\mathbb P^1(\F)$ if and only if
$A_T$ and $B_T$ are linearly dependent and not both zero.  In that case the
contributed projective parameter is the unique point satisfying
\[
        \alpha A_T+\beta B_T=0 .
\]
\end{lemma}

\begin{proof}
A word $Y$ is explained by $C$ on $S=D\setminus T$ if and only if
$Y-c$ is supported on $T$ for some $c\in C$, which is equivalent to
$\Syn(Y)$ lying in the span of the parity-check columns indexed by $T$.  The
standard locator recurrence says that this is equivalent to
\[
        H_{t,j}(\Syn(Y))\boldsymbol\ell_T=0 .
\]
Applying this to $Y=f+zg$ gives $A_T+zB_T=0$.  Under this incidence, simultaneous
explanation of $f$ and $g$ on $S$ is equivalent to $B_T=0$, because
$A_T=(A_T+zB_T)-zB_T$.  This proves the finite-slope statement.  The projective
statement is the same calculation with $Y=\alpha f+\beta g$; if
$A_T=B_T=0$ then both coefficient words are already explained on $S$, so the
point is contained rather than bad.
\end{proof}

Let $\mathcal S_{\mathcal C}^{\ge a}$ be the quotient-remainder support family
from \cref{sec:quotient-support-upper}, and let
\[
        \mathcal T_{\mathcal C}^{\ge a}
        :=\{D\setminus S:S\in\mathcal S_{\mathcal C}^{\ge a}\}
\]
be the corresponding co-support family.  For a received line define
\[
        \Theta_{\mathcal C}^{\rm aff}(f,g;a)
        :=\{\theta_T:T\in\mathcal T_{\mathcal C}^{\ge a},
              \theta_T\text{ is defined by }\cref{lem:exact-line-image-map}\}
\]
and define $\Theta_{\mathcal C}^{\rm proj}(f,g;a)$ analogously using the
projective point of \cref{lem:exact-line-image-map}.

\begin{proposition}[distinct-parameter quotient ledger for lines]
\label{prop:distinct-parameter-line-ledger}
For every received line $f+zg$, the finite affine slopes that are support-wise
MCA-bad and have a witness support in
$\mathcal S_{\mathcal C}^{\ge a}$ are exactly the elements of
$\Theta_{\mathcal C}^{\rm aff}(f,g;a)$.  Hence the exact quotient-branch line
numerator is
\[
        |\Theta_{\mathcal C}^{\rm aff}(f,g;a)|
        \le U_{\mathcal C}^{\rm supp}(a).
\]
The corresponding projective-slope numerator is
\[
        |\Theta_{\mathcal C}^{\rm proj}(f,g;a)|
        \le U_{\mathcal C}^{\rm supp}(a).
\]
The same image ledgers upper-bound no-loss CA parameters whose explaining
support lies in the same quotient family.
\end{proposition}

\begin{proof}
This is only \cref{lem:exact-line-image-map} applied to every co-support in the
family, followed by taking a union of the resulting parameter values.  The CA
claim follows because no-loss CA-badness is stronger than having an explained
slope on the support and forbids simultaneous explanation on any support of size
at least $a$.
\end{proof}

The curve analogue is just as explicit.  For a degree-$d$ power curve
\[
        W_\gamma=f_0+\gamma f_1+\cdots+\gamma^d f_d,
        \qquad \gamma\in\F,
\]
put $u_i=\Syn(f_i)$ and, for $T\subseteq D$, define
\[
        V_{i,T}:=H_{t,j}(u_i)\boldsymbol\ell_T\in\F^t .
\]
Let
\[
        \Gamma_T^{(d)}
        :=\{\gamma\in\F:
             V_{0,T}+\gamma V_{1,T}+\cdots+\gamma^dV_{d,T}=0,
             \ (V_{0,T},\ldots,V_{d,T})\ne0\} .
\]
The nonzero tuple condition is exactly support-wise noncontainment of the
coefficient tuple on the fixed support.

\begin{lemma}[exact support-image map for finite-parameter power curves]
\label{lem:exact-curve-image-map}
For a fixed support $S=D\setminus T$, the set of support-wise noncontained
curve-bad finite parameters is exactly $\Gamma_T^{(d)}$.  In particular
$|\Gamma_T^{(d)}|\le d$.
\end{lemma}

\begin{proof}
The same syndrome-locator recurrence applied to $W_\gamma$ gives the displayed
vector equation.  If every $V_{i,T}$ is zero, then every coefficient word
$f_i$ is already explained on $S$, so the parameter is contained.  Otherwise at
least one scalar coordinate of
$V_{0,T}+\gamma V_{1,T}+\cdots+\gamma^dV_{d,T}$ is a nonzero polynomial in
$\gamma$ of degree at most $d$, so it has at most $d$ roots.
\end{proof}

Define the exact quotient curve image
\[
        \Gamma_{\mathcal C}^{(d)}(f_0,\ldots,f_d;a)
        :=\bigcup_{T\in\mathcal T_{\mathcal C}^{\ge a}}\Gamma_T^{(d)} .
\]

\begin{proposition}[distinct-parameter quotient ledger for curves]
\label{prop:distinct-parameter-curve-ledger}
For a finite-parameter degree-$d$ power curve, the quotient-branch support-wise
curve-MCA numerator is exactly
\[
        |\Gamma_{\mathcal C}^{(d)}(f_0,
        \ldots,f_d;a)| .
\]
It satisfies the safe bounds
\[
        |\Gamma_{\mathcal C}^{(d)}(f_0,\ldots,f_d;a)|
        \le d\,U_{\mathcal C}^{\rm supp}(a)
        \le d\,U_{\mathcal C}^{\rm sum}(a).
\]
The same image ledger upper-bounds no-loss curve-CA parameters whose explaining
support lies in the quotient family.
\end{proposition}

\begin{proof}
Apply \cref{lem:exact-curve-image-map} to each support and take the union of the
resulting parameter sets.  The displayed inequalities follow from
$|\Gamma_T^{(d)}|\le d$ and from the support-union bound of
\cref{prop:divisor-union-support-ledger}.
\end{proof}

\begin{definition}[fixed-support parameter gcd]
\label{def:fixed-support-parameter-gcd}
Let $C=\RS[\F,D,k]$, let $q=|\F|$, and let
\[
        W_\Gamma=f_0+\Gamma f_1+\cdots+\Gamma^d f_d,
        \qquad \Gamma\in\F,
\]
be a finite-parameter degree-$d$ power curve of received words.  Fix a
co-support $T\subseteq D$, put $S=D\setminus T$, and assume
$|S|>k$.  With the notation of \cref{lem:exact-curve-image-map}, write
\[
        P_{T,m}(\Gamma)
        :=\sum_{i=0}^d \Gamma^i (V_{i,T})_m,
        \qquad 0\le m<t,
        \qquad t=|S|-k .
\]
Define the fixed-support contribution polynomial
$g_T^{(d)}(\Gamma)\in\F[\Gamma]$ as follows.  If all coordinate polynomials
$P_{T,m}$ are identically zero, set $g_T^{(d)}=1$.  Otherwise let
$g_T^{(d)}$ be the monic greatest common divisor of the nonzero polynomials
among $P_{T,0},\ldots,P_{T,t-1}$.
\end{definition}

\begin{lemma}[fixed-support gcd equals the noncontained parameter set]
\label{lem:fixed-support-gcd-exact}
In the setting of \cref{def:fixed-support-parameter-gcd}, the finite parameters
$\gamma\in\F$ for which $W_\gamma$ is explained on $S=D\setminus T$ and the
coefficient tuple $(f_0,\ldots,f_d)$ is not simultaneously explained on $S$ are
exactly the roots in $\F$ of $g_T^{(d)}$.  In particular
\[
        \deg g_T^{(d)}\le d,
\]
with equality possible only before taking common factors across different
supports.
\end{lemma}

\begin{proof}
By \cref{lem:exact-curve-image-map}, the fixed support explains $W_\gamma$ exactly
when
\[
        P_{T,m}(\gamma)=0\qquad(0\le m<t).
\]
If every coordinate polynomial is identically zero, then every coefficient
vector $V_{i,T}$ is zero, so every coefficient word $f_i$ is already explained on
$S$; this is the contained branch and contributes no noncontained parameter.
The convention $g_T^{(d)}=1$ gives the empty root set, as required.

Otherwise the common zeros of the coordinate polynomials are exactly the roots
of their monic gcd after zero coordinates are discarded.  Since at least one
coordinate polynomial is nonzero of degree at most $d$, the gcd also has degree
at most $d$.  The nonzero tuple condition in \cref{lem:exact-curve-image-map} is
precisely the assertion that the coordinate polynomials are not all zero.
\end{proof}

\begin{definition}[quotient image lcm polynomial]
\label{def:quotient-image-lcm-polynomial}
Let $\mathcal T$ be any finite family of co-supports $T\subseteq D$ with
$|D\setminus T|>k$.  For a degree-$d$ finite-parameter power curve define
\[
        R_{\mathcal T}^{(d)}(\Gamma)
        :=\operatorname{rad}\operatorname{lcm}_{T\in\mathcal T} g_T^{(d)}(\Gamma),
\]
where factors equal to $1$ are ignored and the empty lcm is $1$.  For the quotient branch at agreement
threshold $a>k$ and divisor set $\mathcal C$, we use
\[
        \mathcal T=\mathcal T_{\mathcal C}^{\ge a}
        :=\{D\setminus S:S\in\mathcal S_{\mathcal C}^{\ge a}\}.
\]
The finite-field root-count polynomial is
\[
        R_{\mathcal T,F}^{(d)}(\Gamma)
        :=\gcd\bigl(R_{\mathcal T}^{(d)}(\Gamma),\Gamma^q-\Gamma\bigr).
\]
Because $R_{\mathcal T}^{(d)}$ is squarefree, the degree of
$R_{\mathcal T,F}^{(d)}$ is the number of $\F$-rational roots of
$R_{\mathcal T}^{(d)}$.
\end{definition}

\begin{theorem}[exact finite-parameter quotient image ledger]
\label{thm:exact-quotient-image-lcm-ledger}
Let $a>k$, let $\mathcal C$ be a finite divisor set, and let
$\mathcal T_{\mathcal C}^{\ge a}$ be the quotient co-support family.  For every
finite-parameter degree-$d$ power curve, the exact quotient-branch support-wise
curve-MCA numerator is
\[
        \boxed{
        \left|\Gamma_{\mathcal C}^{(d)}(f_0,\ldots,f_d;a)\right|
        =\deg R_{\mathcal T_{\mathcal C}^{\ge a},F}^{(d)} .
        }
\]
Consequently
\[
        \left|\Gamma_{\mathcal C}^{(d)}(f_0,\ldots,f_d;a)\right|
        \le \deg R_{\mathcal T_{\mathcal C}^{\ge a}}^{(d)}
        \le d\,U_{\mathcal C}^{\rm supp}(a),
\]
and the final support count may be evaluated by the exact divisor-block formula
of \cref{thm:exact-divisor-block-support-ledger}.  The same lcm ledger
upper-bounds no-loss curve-CA parameters whose explaining support lies in the
same quotient family.
\end{theorem}

\begin{proof}
By \cref{lem:fixed-support-gcd-exact}, for each fixed co-support $T$ the
noncontained parameters witnessed on $D\setminus T$ are exactly the $\F$-roots of
$g_T^{(d)}$.  Taking the union over
$T\in\mathcal T_{\mathcal C}^{\ge a}$ is therefore the same as taking the
$\F$-rational root set of the squarefree lcm
$R_{\mathcal T_{\mathcal C}^{\ge a}}^{(d)}$.  Over a finite field, the
$\F$-rational roots of a squarefree polynomial are counted by the degree of its
gcd with $\Gamma^q-\Gamma$, proving the displayed equality.

The degree bound follows from $\deg g_T^{(d)}\le d$ and from the fact that a
squarefree lcm has degree at most the sum of the degrees of the factors before
coalescing.  The no-loss curve-CA claim follows because a no-loss CA parameter
with an explaining support in the quotient family is, on that support,
noncontained; otherwise the coefficient tuple would have a simultaneous
explanation on a support of size at least $a$.
\end{proof}

\begin{corollary}[affine lines as the degree-one lcm ledger]
\label{cor:affine-line-lcm-ledger}
For an affine received line $f+Zg$, take $d=1$, $f_0=f$, and $f_1=g$.  Then
$R_{\mathcal T_{\mathcal C}^{\ge a}}^{(1)}(Z)$ is the squarefree product of the
linear factors $Z-z$ for the distinct finite quotient-witnessed support-wise
MCA-bad slopes.  Hence
\[
        \#\{\text{finite quotient-witnessed MCA-bad slopes}\}
        =\deg R_{\mathcal T_{\mathcal C}^{\ge a}}^{(1)} .
\]
The same polynomial upper-bounds no-loss CA finite slopes whose explaining
support lies in the declared quotient family.
\end{corollary}

\begin{proof}
For $d=1$, each nonconstant fixed-support gcd has degree at most one.  Its root,
if present, is exactly the unique finite slope supplied by
\cref{lem:exact-line-image-map}.  Taking the squarefree lcm over the declared
co-support family therefore multiplies one copy of $Z-z$ for each distinct
finite slope and no copies for supports that contribute no finite noncontained
slope.
\end{proof}

\begin{definition}[projective line lcm ledger]
\label{def:projective-line-lcm-ledger}
For a projective received line $Z_0f+Z_1g$ and a co-support $T$ with
$|D\setminus T|>k$, form the homogeneous linear coordinate forms
\[
        L_{T,m}(Z_0,Z_1):=Z_0(A_T)_m+Z_1(B_T)_m,
        \qquad 0\le m<t .
\]
If all $L_{T,m}$ are zero, set $h_T^{\rm proj}=1$.  Otherwise let
$h_T^{\rm proj}$ be the gcd of the nonzero linear forms, normalized by a fixed
monomial order.  Thus $h_T^{\rm proj}$ is either $1$ or a homogeneous linear
form.  For a co-support family $\mathcal T$ define
\[
        R_{\mathcal T}^{\rm proj}(Z_0,Z_1)
        :=\operatorname{rad}\operatorname{lcm}_{T\in\mathcal T}h_T^{\rm proj}.
\]
\end{definition}

\begin{corollary}[exact projective quotient image ledger]
\label{cor:projective-line-lcm-ledger}
For projective slopes, the exact quotient-branch support-wise MCA numerator is
\[
        \deg R_{\mathcal T_{\mathcal C}^{\ge a}}^{\rm proj}.
\]
The same degree upper-bounds no-loss projective-CA parameters whose explaining
support lies in the quotient family.
\end{corollary}

\begin{proof}
A nonzero family of homogeneous linear forms on $\mathbb P^1$ has a common zero
if and only if the forms have a common homogeneous linear factor; this factor
cuts out exactly one $\F$-rational projective point.  The all-zero case is the
contained branch and contributes no bad point, by convention.  The squarefree
lcm therefore contains one linear factor for each distinct projective parameter
and no duplicates.
\end{proof}

\begin{definition}[lcm quotient-image certificate]
\label{def:lcm-quotient-image-certificate}
An lcm quotient-image certificate for a line or finite-parameter curve consists
of:
\begin{enumerate}[label=\textup{(\roman*)},leftmargin=2.2em]
\item the divisor set $\mathcal C$, the agreement threshold $a$, and either the
exact support-union certificate of \cref{def:divisor-block-support-ledger-certificate}
or an explicit hash/orbit enumeration of the co-support family
$\mathcal T_{\mathcal C}^{\ge a}$;
\item for every represented co-support $T$, the locator vector
$\boldsymbol\ell_T$ and the syndrome-window coefficient vectors $V_{i,T}$;
\item the coordinate polynomials $P_{T,m}(\Gamma)$, the fixed-support gcd
$g_T^{(d)}$, and a proof that the all-zero case has been assigned to the
contained branch;
\item the squarefree lcm $R_{\mathcal T}^{(d)}$ and, over the finite field, the
root-count factor $R_{\mathcal T,F}^{(d)}=\gcd(R_{\mathcal T}^{(d)},\Gamma^q-\Gamma)$;
\item the declared numerator $\deg R_{\mathcal T,F}^{(d)}$ for finite curves or
$\deg R_{\mathcal T}^{\rm proj}$ for projective lines.
\end{enumerate}
The verifier checks the locator split condition, the syndrome recurrence, the
gcd and lcm arithmetic, squarefreeness, and the final finite-field root count.
This certificate replaces a support count by an exact distinct-parameter count
inside the declared quotient branch.
\end{definition}

\begin{remark}[status after the lcm image milestone]
\label{rem:lcm-image-milestone-status}
The quotient ledgers now separate three tasks that were previously entangled.
First, \cref{thm:quotient-remainder-deep-floor} supplies lower floors by
heaviest generated-field prefix fibers.  Second,
\cref{thm:exact-divisor-block-support-ledger} counts the declared quotient
support union with divisor overlaps handled exactly.  Third,
\cref{thm:exact-quotient-image-lcm-ledger,cor:affine-line-lcm-ledger,cor:projective-line-lcm-ledger}
push that support family through the fixed-locator equations and coalesce
duplicate finite, projective, and curve parameters by gcd/lcm arithmetic.  The
remaining quotient-side obstruction is not duplicate coalescing inside a declared
quotient family; it is the structural safe-side theorem that every bad parameter
in the intended regime either belongs to this quotient image, to the tangent
branch, or to the aperiodic/extension buckets bounded elsewhere.
\end{remark}

\begin{definition}[quotient image certificate]
\label{def:quotient-image-certificate}
A quotient image certificate for a line or curve ledger consists of:
\begin{enumerate}[label=\textup{(\roman*)}]
\item the divisor set $\mathcal C$ and the agreement threshold $a$;
\item a deterministic description, hash, or orbit decomposition of
$\mathcal T_{\mathcal C}^{\ge a}$;
\item for every represented co-support $T$, the locator coefficients
$\boldsymbol\ell_T$ and the syndrome-window vectors $A_T,B_T$ for a line, or
$V_{0,T},\ldots,V_{d,T}$ for a curve;
\item the declared parameter value, parameter bucket, or the declaration that no
parameter is contributed;
\item a bucket table proving the number of distinct parameters after duplicate
coalescing.
\end{enumerate}
The verifier checks the locator split condition, the Hankel equations, the
noncontainment vector, and the bucket consistency.  This certificate is the
preferred replacement for a raw quotient-support count whenever the support count
is much larger than the actual parameter image.
\end{definition}

\begin{remark}[status of the refinement]
\label{rem:distinct-ledger-status}
The propositions above are exact for the declared quotient branch and are a real
refinement over the support ledger: the numerator is an image size, not a support
count.  They still do not prove that all bad parameters are quotient-periodic.
The remaining safe-side theorem for the full prize is the aperiodic Hankel-packing
bound: after tangent and quotient image branches are removed, one must bound the
parameters coming from co-supports outside the declared quotient families.
\end{remark}


\subsection{Aperiodic Hankel chart atlas}
\label{sec:aperiodic-hankel-certificates}

The quotient image ledger of \cref{sec:quotient-distinct-parameter} counts bad
parameters whose witness support lies in a declared quotient-remainder family.
The complementary safe-side problem is to control all other supports.  This
subsection is the contributor-facing version of that problem.  It does not claim
that the needed aperiodic eliminants always exist.  It gives a finite chart
atlas, with precise ideals and pivots, so that a verifier can check a supplied
eliminant and convert it into a bad-parameter upper bound.

The main point is to avoid one huge saturated incidence ideal.  We split the
aperiodic branch into:
\begin{enumerate}[label=\textup{(\roman*)},leftmargin=2.2em]
\item a regular overdetermined Hankel bucket, where a single maximal minor in the
line parameter already gives an eliminant;
\item finite affine pivot charts, one for each noncontainment coordinate;
\item a separate projective-infinity chart, if projective slopes are being
counted;
\item finite-parameter curve pivot charts, one for a nonzero coefficient of one
coordinate polynomial;
\item a singular bucket, where all regular minors vanish or the projection to the
parameter line remains positive-dimensional after the previous splits.
\end{enumerate}
Only the last bucket is genuinely unresolved.  A staircase scanner should keep it
as a named obstruction rather than hiding it inside the word ``aperiodic.''

Fix an exact agreement value $A\ge k$ and put
\[
        j=j_A:=n-A,
        \qquad
        t=t_A:=A-k,
        \qquad
        R=n-k .
\]
Let
\[
        P_D(X):=\prod_{x\in D}(X-x).
\]
A monic degree-$j$ co-support locator is written
\[
        L_{\ell}(X)=X^j+\ell_{j-1}X^{j-1}+\cdots+\ell_0,
        \qquad \ell=(\ell_0,\ldots,\ell_{j-1})\in\mathbb A^j .
\]
Let $R_{D,j}(\ell;X)$ be the remainder of $P_D(X)$ modulo $L_{\ell}(X)$, and let
$I_{\rm split,j}$ be the ideal generated by the $j$ coefficients of
$R_{D,j}$.  Since $P_D$ is squarefree, the $F$-points of
$V(I_{\rm split,j})$ are exactly the squarefree $D$-split monic degree-$j$
locators.

For a word $y:D\to F$ write $\Syn(y)\in F^R$ for the parity-check syndrome used
above.  Let
\[
        \widetilde\ell=(\ell_0,\ldots,\ell_{j-1},1)^t .
\]
For an affine line $f+zg$, put $u=\Syn(f)$ and $v=\Syn(g)$ and define
\[
        A(\ell):=H_{t,j}(u)\widetilde\ell,
        \qquad
        B(\ell):=H_{t,j}(v)\widetilde\ell
        \qquad\in F^t .
\]
For a fixed split locator, the support-wise finite-slope condition is exactly
\[
        A(\ell)+ZB(\ell)=0,
        \qquad
        B(\ell)\ne0 .
\]
The inequality is the noncontainment condition.

\subsubsection{Regular overdetermined bucket}

Assume first that
\[
        t_A\ge j_A+1,
        \qquad\text{equivalently}\qquad
        2A\ge n+k+1 .
\]
Set
\[
        M_A(Z):=H_{t_A,j_A}(u)+Z H_{t_A,j_A}(v).
\]
This is a $t_A\times(j_A+1)$ matrix with entries affine-linear in $Z$.

\begin{definition}[regular Hankel minor certificate]
\label{def:hankel-regularity-certificate}
A regular Hankel minor certificate at exact agreement $A$ is a row set
$I_A\subseteq\{0,\ldots,t_A-1\}$ of size $j_A+1$ such that
\[
        \Delta_A(Z)
        :=\det M_A(Z)_{I_A,\{0,\ldots,j_A\}}
\]
is not the zero polynomial in $F[Z]$.  If no such row set exists, the line is
called \emph{Hankel-singular at $A$}.
\end{definition}

\begin{lemma}[regular exact-agreement eliminant]
\label{lem:regular-exact-agreement-eliminant}
Assume $t_A\ge j_A+1$, and suppose $\Delta_A(Z)$ is a regular Hankel minor
certificate.  Every finite support-wise MCA-bad slope with a witness support of
exact size $A$ is a root of $\Delta_A$.  Hence the number of such slopes is at
most
\[
        \deg\Delta_A\le j_A+1=n-A+1 .
\]
The same root containment holds for no-loss CA-bad slopes with exact agreement
$A$.
\end{lemma}

\begin{proof}
Let $z$ be a bad finite slope with witness support $S$ of size $A$, and put
$T=D\setminus S$.  Then $|T|=j_A$, and the locator vector
$\widetilde\ell_T$ is a nonzero kernel vector of $M_A(z)$.  Therefore the column
rank of $M_A(z)$ is less than $j_A+1$, so every $(j_A+1)\times(j_A+1)$ maximal
minor vanishes at $z$, including $\Delta_A(z)$.  The degree bound follows because
$\Delta_A$ is the determinant of a $(j_A+1)\times(j_A+1)$ matrix whose entries
are affine-linear in $Z$.
\end{proof}

\begin{theorem}[regular closed-range Hankel packing]
\label{thm:regular-closed-ball-hankel-packing}
Let $a\ge k$ and suppose that for each exact agreement
$A\in\{a,a+1,\ldots,n\}$ with $t_A\ge j_A+1$ a regular Hankel minor certificate
$\Delta_A$ is supplied.  Then all bad finite slopes whose exact witness size is
in this regular bucket lie in the root set of
\[
        \Delta_{\ge a}(Z):=\prod_{A=a}^{n}\Delta_A(Z),
\]
with factors omitted for exact agreements outside the overdetermined range.
Consequently their number is at most
\[
        \sum_A \deg\Delta_A .
\]
If exact root sets are supplied, roots already charged to the tangent or
quotient-image ledgers may be removed from this count by duplicate coalescing.
\end{theorem}

\begin{proof}
Apply \cref{lem:regular-exact-agreement-eliminant} at the exact size of the
witness support and take the union of the resulting root sets.  The degree sum is
a safe union bound.
\end{proof}

\paragraph{Canonical regular rank-drop gcds.}
The regular-minor certificate below can be made canonical.  Instead of choosing
one nonzero maximal minor, take the greatest common divisor of all nonzero
maximal minors.  This gives exactly the finite-parameter values at which the
Hankel pencil can drop rank, before the split-locator and noncontainment tests
are imposed.  Thus it is an upper ledger for the regular bucket, and often a much
smaller one than the degree of a single selected minor.

Fix an exact agreement value $A\ge k$, and keep the notation
\[
        j=j_A:=n-A,
        \qquad
        t=t_A:=A-k,
        \qquad
        r=n-k .
\]
The regular overdetermined range is
\[
        t_A>j_A,
        \qquad\text{i.e.}\qquad
        t_A\ge j_A+1,
        \qquad\text{equivalently}\qquad
        2A\ge n+k+1 .
\]

\begin{definition}[canonical affine rank-drop gcd]
\label{def:canonical-affine-rankdrop-gcd}
Fix an affine received line $f+Zg$, put
$u=\Syn(f)$ and $v=\Syn(g)$, and assume $t_A\ge j_A+1$.  Set
\[
        M_A(Z):=H_{t_A,j_A}(u)+Z H_{t_A,j_A}(v).
\]
Let $\mathfrak M_A(Z)$ be the finite set of all $(j_A+1)\times(j_A+1)$ maximal
minors of $M_A(Z)$.  If every element of $\mathfrak M_A(Z)$ is the zero
polynomial, the exact-agreement bucket is called \emph{rank-drop singular}.
Otherwise define
\[
        G_A^{\rm aff}(Z)
        :=\operatorname{rad}\gcd\{\Delta(Z):
              \Delta\in\mathfrak M_A(Z),\ \Delta\ne0\},
\]
with the gcd normalized to be monic.  Over the finite field $\F$ of size $q$ set
\[
        G_{A,F}^{\rm aff}(Z):=
        \gcd\bigl(G_A^{\rm aff}(Z),Z^q-Z\bigr).
\]
Thus $G_{A,F}^{\rm aff}$ is the squarefree polynomial whose roots are precisely
those $F$-rational parameters at which all nonzero maximal minors vanish.
\end{definition}

\begin{theorem}[canonical exact-agreement rank-drop ledger]
\label{thm:canonical-affine-rankdrop-ledger}
Assume $t_A\ge j_A+1$ and the exact-agreement bucket is not rank-drop singular.
Every finite support-wise MCA-bad slope with an exact witness support of size
$A$ is a root of $G_A^{\rm aff}$.  Consequently the number of such slopes is at
most
\[
        \deg G_{A,F}^{\rm aff}
        \le \deg G_A^{\rm aff}
        \le j_A+1=n-A+1 .
\]
The same bound holds for no-loss CA-bad slopes with exact agreement $A$.
\end{theorem}

\begin{proof}
Let $z\in F$ be support-wise bad with exact witness support $S$ of size $A$, and
put $T=D\setminus S$.  By \cref{lem:support-locator-syndrome-recurrence}, the
locator vector $\boldsymbol\ell_T$ is a nonzero kernel vector of $M_A(z)$.
Therefore $\operatorname{rank}M_A(z)<j_A+1$, so every maximal minor of $M_A(Z)$
vanishes at $Z=z$.  In particular every nonzero maximal minor vanishes at $z$,
and hence their gcd $G_A^{\rm aff}$ vanishes at $z$.

The finite-field root count is exactly $\deg G_{A,F}^{\rm aff}$ because
$G_{A,F}^{\rm aff}$ is squarefree and divides $Z^q-Z$.  Each maximal minor has
degree at most $j_A+1$, since it is the determinant of a matrix whose entries are
affine-linear in $Z$; the gcd cannot have larger degree.  The no-loss CA claim is
identical, because no-loss CA-badness with exact agreement $A$ also supplies an
explaining support of size $A$ and hence the same locator-kernel condition.
\end{proof}

\begin{remark}[why this improves the one-minor certificate]
\label{rem:gcd-rankdrop-improves-one-minor}
\Cref{lem:regular-exact-agreement-eliminant} used one certified nonzero maximal
minor and therefore counted the roots of a single determinant.  The polynomial
$G_A^{\rm aff}$ uses all maximal minors and counts their common finite-field root
set.  It is never worse than any one-minor certificate and can be much smaller;
it is also canonical, so two scanners that use the same syndrome convention print
the same regular-bucket polynomial.
\end{remark}

\begin{definition}[closed-ball rank-drop lcm]
\label{def:closed-ball-rankdrop-lcm}
For an agreement threshold $a>k$, let
\[
        \mathcal A_{\rm reg}(a):=
        \{A\in\{a,a+1,\ldots,n\}: t_A\ge j_A+1
          \text{ and the bucket }A\text{ is not rank-drop singular}\}.
\]
Define
\[
        G_{\ge a}^{\rm aff}(Z)
        :=\operatorname{rad}\operatorname{lcm}_{A\in\mathcal A_{\rm reg}(a)}
          G_A^{\rm aff}(Z),
        \qquad
        G_{\ge a,F}^{\rm aff}(Z):=
        \gcd\bigl(G_{\ge a}^{\rm aff}(Z),Z^q-Z\bigr),
\]
with the empty lcm equal to $1$.  Exact agreement values
$A\in\{a,\ldots,n\}\setminus\mathcal A_{\rm reg}(a)$ are recorded as residual
buckets: they are either underdetermined ($t_A<j_A+1$) or rank-drop singular.
\end{definition}

\begin{corollary}[canonical regular closed-ball Hankel packing]
\label{cor:canonical-regular-closed-ball-hankel-packing}
All finite support-wise MCA-bad slopes whose exact witness size lies in
$\mathcal A_{\rm reg}(a)$ are roots of $G_{\ge a}^{\rm aff}$.  Hence their number
is at most
\[
        \deg G_{\ge a,F}^{\rm aff}
        \le \sum_{A\in\mathcal A_{\rm reg}(a)}(n-A+1).
\]
The same bound holds for no-loss CA-bad slopes in the same regular buckets.
\end{corollary}

\begin{proof}
Apply \cref{thm:canonical-affine-rankdrop-ledger} at the exact agreement size of
the witness and take the union over $A$.  The squarefree lcm coalesces duplicate
roots across exact-agreement buckets; the final gcd with $Z^q-Z$ counts only
finite-field roots.  The displayed degree sum is the safe bound obtained before
coalescing.
\end{proof}

\begin{definition}[paid-root removal for the regular branch]
\label{def:paid-root-removal-regular-branch}
Suppose a squarefree polynomial $P_{\rm paid,F}(Z)$ is supplied whose roots are
exactly the finite-field parameters already charged to tangent/common-line and
quotient-image ledgers.  Define the regular residual polynomial
\[
        G_{{\rm reg}\setminus{\rm paid},F}^{\rm aff}(Z)
        :=\frac{G_{\ge a,F}^{\rm aff}(Z)}
        {\gcd(G_{\ge a,F}^{\rm aff}(Z),P_{\rm paid,F}(Z))} .
\]
Its degree is the number of regular rank-drop parameters not already paid by the
previous ledgers.
\end{definition}

\begin{corollary}[quotient/tangent-removed regular aperiodic numerator]
\label{cor:quotient-tangent-removed-regular-numerator}
If $P_{\rm paid,F}$ is exact in the sense of
\cref{def:paid-root-removal-regular-branch}, then the number of finite slopes in
the regular overdetermined branch that remain after removing the paid
quotient/tangent parameters is at most
\[
        \deg G_{{\rm reg}\setminus{\rm paid},F}^{\rm aff} .
\]
This is a safe aperiodic numerator for the regular buckets.  The underdetermined
and rank-drop singular exact-agreement buckets must still be handled by pivot
charts or by a separately named residual obstruction.
\end{corollary}

\begin{proof}
By \cref{cor:canonical-regular-closed-ball-hankel-packing}, every regular-bucket
bad slope is a root of $G_{\ge a,F}^{\rm aff}$.  Removing the roots of the exact
paid polynomial deletes precisely those already charged parameters from this
finite root set.  The remaining root set is counted by the displayed quotient
because all polynomials are squarefree divisors of $Z^q-Z$.
\end{proof}

\begin{definition}[finite-parameter curve rank-drop gcd]
\label{def:curve-rankdrop-gcd}
For a finite-parameter degree-$d$ curve
\[
        W_\Gamma=f_0+\Gamma f_1+\cdots+\Gamma^d f_d,
\]
put $u_i=\Syn(f_i)$ and, at exact agreement $A$, define
\[
        M_A^{(d)}(\Gamma):=
        \sum_{i=0}^{d}\Gamma^i H_{t_A,j_A}(u_i).
\]
When $t_A\ge j_A+1$ and not all maximal minors of $M_A^{(d)}$ vanish
identically, set
\[
        G_A^{(d)}(\Gamma)
        :=\operatorname{rad}\gcd\{\Delta(\Gamma):
              \Delta\text{ is a nonzero maximal minor of }M_A^{(d)}\},
\]
and
\[
        G_{A,F}^{(d)}(\Gamma):=
        \gcd\bigl(G_A^{(d)}(\Gamma),\Gamma^q-\Gamma\bigr).
\]
\end{definition}

\begin{theorem}[canonical curve rank-drop ledger]
\label{thm:canonical-curve-rankdrop-ledger}
In a non-singular overdetermined exact-agreement bucket, every finite
curve-MCA-bad parameter is a root of $G_A^{(d)}$.  Consequently the number of
such parameters is at most
\[
        \deg G_{A,F}^{(d)}
        \le \deg G_A^{(d)}
        \le d(j_A+1).
\]
The same bound holds for no-loss curve-CA parameters in the same bucket.  The
closed-ball version is obtained by taking the squarefree lcm of the polynomials
$G_A^{(d)}$ over the regular exact-agreement buckets and then taking the gcd with
$\Gamma^q-\Gamma$.
\end{theorem}

\begin{proof}
A curve-bad parameter with exact witness support $S$ gives, for
$T=D\setminus S$, a nonzero locator vector in the kernel of
$M_A^{(d)}(\gamma)$.  Hence every maximal minor vanishes at $\gamma$, and the
same gcd argument as in \cref{thm:canonical-affine-rankdrop-ledger} applies.
Each matrix entry has degree at most $d$ in $\Gamma$, so each maximal minor has
degree at most $d(j_A+1)$.
\end{proof}

\begin{definition}[regular rank-drop certificate]
\label{def:regular-rankdrop-certificate}
A regular rank-drop certificate for a row, a received affine line or finite
curve, and an agreement threshold $a$ consists of:
\begin{enumerate}[label=\textup{(\roman*)},leftmargin=2.2em]
\item the syndrome vectors of the coefficient words and the exact-agreement range
$A\in\{a,\ldots,n\}$;
\item for every overdetermined exact agreement $A$, either a declaration that all
maximal minors vanish identically, or the canonical gcd polynomial
$G_A^{\rm aff}$, respectively $G_A^{(d)}$, together with ideal-membership or raw
minor arithmetic proving that it is the gcd of all nonzero maximal minors;
\item the squarefree lcm over all non-singular regular buckets and the finite-root
factor obtained by gcd with $Z^q-Z$ or $\Gamma^q-\Gamma$;
\item if paid branches are removed, the exact paid-root polynomial and the finite
root-set quotient of \cref{def:paid-root-removal-regular-branch};
\item a residual-bucket table listing all underdetermined and rank-drop singular
exact agreements that are not covered by the regular gcd ledger.
\end{enumerate}
The verifier checks the syndrome convention, the Hankel matrices, the maximal
minor gcds, squarefree lcms, finite-field root factors, paid-root removal, and the
residual-bucket table.
\end{definition}

\begin{remark}[status after the rank-drop milestone]
\label{rem:rankdrop-milestone-status}
The regular overdetermined part of the aperiodic Hankel problem is now canonical:
one no longer has to choose a minor or accept a degree sum when the common root
set of all maximal minors is smaller.  This still does not solve the near-capacity
underdetermined buckets, which are the hard part of T2 in the prize program.
Those buckets remain assigned to the affine pivot charts, projective-infinity
charts, curve coefficient pivots, or named singular residual obstructions.
\end{remark}

\subsubsection{Finite affine pivot charts}

The regular minor test is intentionally coarse and only works directly in the
overdetermined range.  The next atlas is the one contributors should use to
construct actual aperiodic eliminants.

\begin{definition}[aperiodic locator chart]
\label{def:aperiodic-locator-chart}
An aperiodic locator chart at exact agreement $A$ is a pair
\[
        X=(E_X,\Delta_X),
\]
where $E_X$ is a finite list of polynomial equations in the locator coefficients
$\ell$, and $\Delta_X(\ell)$ is a product of declared nonzero tests.  The chart
represents the constructible set
\[
        V(I_{\rm split,j}+\langle E_X\rangle)\setminus V(\Delta_X).
\]
The equations and inequations must include, or reference, the removal of already
paid tangent/common-code-line and quotient-image branches.  A family of charts
$\mathfrak X_A$ is a valid aperiodic chart cover if the union of their
constructible sets contains every remaining split locator of exact co-support
size $j_A$ not assigned to those earlier ledgers.
\end{definition}

For a coordinate $h\in\{0,\ldots,t-1\}$ define the affine pivot consistency ideal
\[
\begin{aligned}
        J^{\rm aff}_{X,h}:={}&
        I_{\rm split,j}+\langle E_X\rangle
        +\langle A_m(\ell)B_h(\ell)-A_h(\ell)B_m(\ell):0\le m<t\rangle \\
        &+\langle ZB_h(\ell)+A_h(\ell)\rangle
        \subseteq F[\ell_0,\ldots,\ell_{j-1},Z],
\end{aligned}
\]
and saturate it by the pivot denominator:
\[
        \widehat J^{\rm aff}_{X,h}:=
        J^{\rm aff}_{X,h}:(\Delta_X B_h)^\infty .
\]
The equations $A_mB_h-A_hB_m=0$ force $A$ and $B$ to be collinear on the pivot
chart, and $ZB_h+A_h=0$ records the unique finite slope.

\begin{theorem}[pivoted affine-line eliminant certificate]
\label{thm:aperiodic-affine-eliminant}
Let $X$ be an aperiodic locator chart and let $h$ be a noncontainment pivot.  If
there is a nonzero polynomial
\[
        Q_{X,h}(Z)\in \widehat J^{\rm aff}_{X,h}\cap F[Z]
\]
of degree $D_{X,h}$, then the number of finite support-wise MCA-bad slopes in
that chart with $B_h\ne0$ is at most $\min\{|F|,D_{X,h}\}$.  The same bound holds
for no-loss CA-bad slopes in that pivot chart.
\end{theorem}

\begin{proof}
A bad slope in the chart has a split locator with $\Delta_X\ne0$ and, on this
pivot chart, $B_h\ne0$.  The Hankel incidence equation gives
$A+ZB=0$, hence the point $(\ell,Z)$ satisfies all generators of
$J^{\rm aff}_{X,h}$ and survives the saturation.  Therefore $Q_{X,h}(Z)=0$.
A nonzero degree-$D_{X,h}$ polynomial over $F$ has at most $D_{X,h}$ roots.
\end{proof}

\begin{remark}[why the pivot split is useful]
\label{rem:pivot-split-useful}
The saturation by the whole vector $B$ in the unsplit incidence ideal is replaced
by the explicit open cover $B_h\ne0$.  A scanner can try all $h$, produce small
univariate eliminants on successful pivots, and report the residual locus
$B_0=\cdots=B_{t-1}=0$ as contained for support-wise MCA.  For no-loss CA, that
residual locus also contributes no bad slope at the same support, because if
$A=B=0$ then both $f$ and $g$ are explained on the witness support.
\end{remark}

\subsubsection{Projective-slope split}

For projective parameters $[Z_0:Z_1]$ the incidence equation is
\[
        Z_0A(\ell)+Z_1B(\ell)=0 .
\]
The finite patch $Z_0\ne0$ is exactly the affine pivot atlas above with
$Z=Z_1/Z_0$.  The extra point at infinity is $[0:1]$.

\begin{definition}[projective infinity chart]
\label{def:projective-infinity-chart}
For a chart $X$, the projective-infinity chart is the constructible locus
\[
        I_{\rm split,j}+\langle E_X\rangle+
        \langle B_0,\ldots,B_{t-1}\rangle,
        \qquad
        \Delta_X\ne0,
        \quad A\ne0 .
\]
If it is nonempty, it contributes the single projective parameter $[0:1]$.
If it is empty, the infinity point contributes nothing on $X$.
\end{definition}

\begin{corollary}[projective eliminant certificate]
\label{cor:projective-eliminant-certificate}
For a chart cover $\mathfrak X_A$, the projective-slope numerator is bounded by
the union of the finite-patch root sets from
\cref{thm:aperiodic-affine-eliminant} together with at most one extra point
$[0:1]$, provided the infinity charts are certified either empty or nonempty.
Equivalently, a homogeneous nonzero eliminant
\[
        Q_X(Z_0,Z_1)
\]
of degree $D_X$ bounds the number of projective bad parameters in that chart by
$D_X$.
\end{corollary}

\begin{proof}
The finite patch is the affine theorem.  The complement of that patch is the
single point $[0:1]$, which is governed by \cref{def:projective-infinity-chart}.
For the homogeneous formulation, every bad projective parameter is a zero of the
homogeneous eliminant, and a nonzero homogeneous polynomial of degree $D_X$ on
$\mathbb P^1(F)$ has at most $D_X$ zeros.
\end{proof}

\subsubsection{Finite-parameter curve charts}

For a degree-$d$ finite power curve
\[
        W_\Gamma=f_0+\Gamma f_1+\cdots+\Gamma^d f_d,
\]
put
\[
        V_i(\ell):=H_{t,j}(\Syn(f_i))\widetilde\ell\in F^t,
        \qquad 0\le i\le d .
\]
The curve incidence equation is
\[
        V_0(\ell)+\Gamma V_1(\ell)+\cdots+\Gamma^d V_d(\ell)=0 .
\]
Curve-MCA noncontainment means that not all coefficient vectors $V_i(\ell)$ are
zero.  Thus the noncontained branch is covered by coefficient pivots
$(i,h)$ with $(V_i)_h\ne0$.

For a chart $X$ and a coefficient pivot $(i,h)$, define
\[
\begin{aligned}
        J^{\rm curve,d}_{X,i,h}:={}&
        I_{\rm split,j}+\langle E_X\rangle \\
        &+\left\langle
          \sum_{r=0}^d \Gamma^r (V_r(\ell))_m:
          0\le m<t
          \right\rangle
        \subseteq F[\ell_0,\ldots,\ell_{j-1},\Gamma],
\end{aligned}
\]
and set
\[
        \widehat J^{\rm curve,d}_{X,i,h}
        :=J^{\rm curve,d}_{X,i,h}:(\Delta_X (V_i)_h)^\infty .
\]

\begin{theorem}[pivoted curve eliminant certificate]
\label{thm:aperiodic-curve-eliminant}
If a chart $X$ and coefficient pivot $(i,h)$ have a nonzero eliminant
\[
        Q_{X,i,h}(\Gamma)\in
        \widehat J^{\rm curve,d}_{X,i,h}\cap F[\Gamma]
\]
of degree $D_{X,i,h}$, then the number of finite degree-$d$ curve-MCA bad
parameters in that chart and pivot is at most $\min\{|F|,D_{X,i,h}\}$.  The same
bound applies to no-loss curve-CA parameters in the same pivot chart.
\end{theorem}

\begin{proof}
Every bad curve parameter on the pivot chart gives a split locator satisfying all
curve-incidence equations and the nonzero pivot condition $(V_i)_h\ne0$.  It
therefore survives the saturation and is a root of the supplied eliminant.
\end{proof}

\subsubsection{Fallback and final certificate}

The eliminant certificates above are strongest when the projection to the
parameter line is finite.  A failed eliminant search should not be discarded; it
should be split further or recorded as a singular bucket.  The following fallback
is safe but usually too weak for a final $2^{-128}$ threshold.

\begin{proposition}[dimension-degree fallback]
\label{prop:dimension-degree-fallback}
Let $Y\subseteq\mathbb A^N_F$ be an affine variety of degree at most $\Delta$ and
Krull dimension at most $e$.  Then
\[
        |Y(F)|\le \Delta |F|^e .
\]
Consequently, if a saturated Hankel incidence chart for line or curve parameters
has degree $\Delta_X$ and dimension $e_X$, then the number of bad parameters in
that chart is at most $\Delta_X|F|^{e_X}$.
\end{proposition}

\begin{proof}
This is the standard affine degree point-count bound, proved by induction on
dimension using generic affine hyperplane sections.  The parameter image is no
larger than the incidence set.
\end{proof}

For a threshold statement, use finitely many charts and agreement levels.  Let
$\mathfrak X_A$ be a chart cover of the declared aperiodic branch at exact
agreement $A$.  If each finite affine pivot chart has a univariate eliminant, set
\[
        B_{\rm ap}^{\rm aff,elim}(a)
        :=\sum_{A=a}^n\sum_{X\in\mathfrak X_A}
          \sum_{h=0}^{t_A-1}\deg Q_{X,h} .
\]
For projective slopes add the certified infinity contribution, and for degree-$d$
curves replace the inner sum by the coefficient-pivot sum
$\sum_{i=0}^d\sum_h\deg Q_{X,i,h}$.  Exact root-union tables may replace these
safe sums whenever duplicate parameters across charts are certified.

\begin{definition}[aperiodic Hankel-packing certificate]
\label{def:aperiodic-hankel-certificate}
An aperiodic Hankel-packing certificate for a row, a line or curve family, and an
agreement threshold $a$ consists of:
\begin{enumerate}[label=\textup{(\roman*)},leftmargin=2.2em]
\item the tangent/common-code-line and quotient-image ledgers that have already
been removed;
\item for each exact agreement $A\in\{a,\ldots,n\}$, either a regular minor
certificate $\Delta_A$ or a declaration that the exact-agreement bucket is
Hankel-singular;
\item for every singular or near-capacity bucket, a finite constructible locator
chart cover $\mathfrak X_A$ as in \cref{def:aperiodic-locator-chart};
\item for affine lines, a list of pivot eliminants $Q_{X,h}(Z)$, or a failed-pivot
record explaining which residual locus remains unsplit;
\item for projective slopes, the finite-patch eliminants plus the infinity-chart
empty/nonempty certificate;
\item for curves, coefficient-pivot eliminants $Q_{X,i,h}(\Gamma)$;
\item a root-union table or a safe degree-sum table producing the declared
aperiodic numerator $B_{\rm ap}(a)$;
\item if any chart uses the dimension-degree fallback, its degree, dimension, and
an explicit statement that the resulting bound is only a fallback, not a sharp
threshold certificate.
\end{enumerate}
The verifier checks locator divisibility by $P_D$, chart coverage of the declared
aperiodic complement, nonzero pivot tests, ideal membership of eliminants after
saturation, and final root-count arithmetic.
\end{definition}

\begin{remark}[contributor workflow]
\label{rem:aperiodic-contributor-workflow}
For each row and exact agreement $A$, contributors should try the following in
order.  First test the overdetermined regular minors.  If a nonzero minor is
found, record its root set and remove roots already paid by tangent or quotient
ledgers.  If all maximal minors vanish identically, move the instance to the
singular bucket.  On the singular or near-capacity bucket, build locator charts,
then split by affine pivots $B_h\ne0$; for projective samplers also test the
infinity chart $B=0$, $A\ne0$; for curves split by coefficient pivots
$(V_i)_h\ne0$.  A chart is complete only when it has either a nonzero eliminant,
a certified empty locus, or a named residual obstruction.  Residual obstructions
should be labelled as quotient-periodic, tangent/common-code-line,
subfield/extension exceptional, or candidate new obstruction.
\end{remark}

\begin{remark}[status]
\label{rem:aperiodic-certificate-status}
\Cref{thm:regular-closed-ball-hankel-packing,thm:aperiodic-affine-eliminant,cor:projective-eliminant-certificate,thm:aperiodic-curve-eliminant}
are proved certificate theorems.  They do not prove that all rows admit small
aperiodic eliminants.  They make the remaining MCA safe-side task concrete: after
tangent and quotient branches are removed, construct small eliminants for the
pivot charts, or identify the singular chart that prevents such an eliminant.
\end{remark}



\begin{lemma}[interleaving transfer]\label{lem:inter}
For every linear code $C\subseteq\F^n$, every $s\ge1$ and every $\delta\in(0,1)$:
\[
\eca\bigl(C^{\equiv s},\delta\bigr)\ \ge\ \eca(C,\delta),
\qquad
\emca\bigl(C^{\equiv s},\delta\bigr)\ \ge\ \emca(C,\delta).
\]
\end{lemma}

\begin{proof}
For $h\in\F^n$ write $h^{\Delta}:=(h,\dots,h)\in(\F^n)^s$. For any codeword tuple $\mathbf c=(c^{(1)},\dots,c^{(s)})\in C^{\equiv s}$, the columns where $h^\Delta$ and $\mathbf c$ agree are $\bigcap_i\{j:c^{(i)}_j=h_j\}\subseteq\{j:c^{(1)}_j=h_j\}$, so $\dist_s(h^\Delta,C^{\equiv s})\ge\dist(h,C)$; taking $c^{(1)}=\dots=c^{(s)}$ gives equality:
$\dist_s(h^\Delta,C^{\equiv s})=\dist(h,C)$.

Fix $(f_1,f_2)$ attaining $\eca(C,\delta)$ (resp.\ $\emca$) and consider the pair $(f_1^\Delta,f_2^\Delta)$ for $C^{\equiv s}$, with the slope $\gamma$ acting row-wise, so $f_1^\Delta+\gamma f_2^\Delta=(f_1+\gamma f_2)^\Delta$. By the displayed equality, $\dist_s\bigl((f_1+\gamma f_2)^\Delta,C^{\equiv s}\bigr)\le\delta$ iff $\dist(f_1+\gamma f_2,C)\le\delta$. For the correlated condition, the pair $(f_1^\Delta,f_2^\Delta)$ viewed as a $2s$-row interleaved word has rows $f_1$ ($s$ times) and $f_2$ ($s$ times); by the same column argument its distance to $(C^{\equiv s})^{\equiv2}=C^{\equiv2s}$ equals $\dist_2((f_1,f_2),C^{\equiv2})$. For the MCA condition, a common set $S$ explains $f_1^\Delta$ and $f_2^\Delta$ iff it explains $f_1$ and $f_2$. Hence $\gamma$ is CA-bad (resp.\ MCA-bad) for $(f_1^\Delta,f_2^\Delta)$ iff it is so for $(f_1,f_2)$, and the maxima over pairs can only grow.
\end{proof}

\section{The main theorem}\label{sec:main}

\begin{theorem}[universal cap]\label{thm:main}
Let $\B\subseteq\F$ be finite fields, let $q=|\F|$, let $D\subseteq\B^\times$ be a multiplicative coset of order $n$, and let $k$ have $\rho=k/n\in(0,1)$; set $C:=\RS[\F,D,k]$. Let $N\mid n$ satisfy $a:=n/N\mid k$ and $(1-\rho)N\ge3$. If
\begin{equation}\label{eq:hyp}
        \binom{N}{\rho N+2}
        \ge
        |\B|\Bigl(\frac qk+1\Bigr),
\end{equation}
then, writing
\[
        \delta_N:=1-\rho-\frac2N,
\]
we have
\[
\eca(C,\delta)
>
\frac1{2k}\Bigl(1-\frac nq\Bigr)
\]
for every integer-admissible radius
\[
        \delta_N\le \delta \le 1-\rho-\frac1n .
\]
In particular,
\[
        \emca(C,\delta)
        >
        \frac1{2k}\Bigl(1-\frac nq\Bigr)
\]
for every
\[
        \delta_N\le \delta < 1-\rho .
\]
Consequently,
\[
        \delta^*_C(\varepsilon^*)
        \le
        1-\rho-\frac2N
\]
for every
\[
        \varepsilon^*
        \le
        \frac1{2k}\Bigl(1-\frac nq\Bigr).
\]
\end{theorem}

\begin{proof}
Note first that $\rho N=k/a\in\mathbb Z$, $\ell_2=\rho N+2\le N-1$ by $(1-\rho)N\ge3$, and $\delta_N=1-\rho-2/N>0$ for the same reason. Also $k+1\le n$.

Fix
\[
        \delta\in[\delta_N,\,1-\rho-\tfrac1n].
\]
Then
\[
        \lfloor \delta n\rfloor \le n-k-1,
\]
so \cref{thm:A} is applicable.

Suppose toward a contradiction that
\[
        \eca(C,\delta)
        \le
        \frac1{2k}\Bigl(1-\frac nq\Bigr).
\]
Apply \cref{thm:A} with $\eta=\frac12$. Its hypothesis holds, so
\[
        \Lst(C^+,\delta)
        \le
        \Bigl\lceil 2q\cdot \eca(C,\delta)\Bigr\rceil
        \le
        \Bigl\lceil\frac{q-n}{k}\Bigr\rceil
        <
        \frac{q-n}{k}+1
        \le
        \frac qk+1 .
\]
On the other hand, by \cref{lem:fiber}(ii) and monotonicity of lists in the radius, since $\delta\ge\delta_N$,
\[
        \Lst(C^+,\delta)
        \ge
        \bigl|\Lst(C^+,\delta_N,u_z)\bigr|
        \ge
        \binom{N}{\ell_2}\Big/|\B|
        \ge
        \frac qk+1
\]
by \eqref{eq:hyp}. This is a contradiction. Hence
\[
        \eca(C,\delta)
        >
        \frac1{2k}\Bigl(1-\frac nq\Bigr)
\]
for every $\delta\in[\delta_N,1-\rho-1/n]$.

In particular, the above holds at $\delta=\delta_N$. By \cref{fact:chain},
\[
        \emca(C,\delta_N)
        \ge
        \eca(C,\delta_N)
        >
        \frac1{2k}\Bigl(1-\frac nq\Bigr).
\]
By support-wise MCA monotonicity, \cref{lem:mca-monotone}, the same lower bound for $\emca(C,\delta)$ holds for every
\[
        \delta\in[\delta_N,1-\rho).
\]
Thus no sub-capacity radius
\[
        \delta\ge 1-\rho-\frac2N
\]
is $\varepsilon^*$-admissible whenever
\[
        \varepsilon^*
        \le
        \frac1{2k}\Bigl(1-\frac nq\Bigr),
\]
which proves the claimed bound on $\delta^*_C(\varepsilon^*)$.
\end{proof}

\begin{remark}[constants]\label{rem:eta}
Choosing $\eta=1-\theta$ in \cref{thm:A} improves the conclusion to $\eca>\frac{1-\theta}k(1-\frac nq)$ at the price of strengthening \eqref{eq:hyp} to $\binom N{\rho N+2}\ge|\B|\bigl(\frac q{\theta k}+1\bigr)$; given the hundreds of bits of slack in \eqref{eq:hyp} throughout the challenge envelope (\cref{tab:numerics}), the error constant is effectively $(1-o(1))/k$. We fix $\eta=\frac12$ for cleanliness.
\end{remark}

\begin{remark}[scope]\label{rem:scope}
No smoothness is assumed beyond the existence of a single divisor $N\mid n$ with $a=n/N\mid k$ and \eqref{eq:hyp}. The theorem applies verbatim to $2$-adic, mixed-radix, or any other FFT-friendly multiplicative domain, and --- via the subfield pigeonhole --- to any extension field $\F\supseteq\B$ with $D\subseteq\B$.

The deep-point conversion is an integer-radius theorem with the condition $\lfloor\delta n\rfloor\le n-k-1$. Therefore the no-loss CA lower bound is stated for
\[
        \delta \le 1-\rho-\frac1n .
\]
The MCA threshold cap is stronger in the range needed for the Proximity Prize: once the support-wise MCA lower bound is proved at
\[
        \delta_N=1-\rho-\frac2N,
\]
monotonicity of support-wise MCA extends it to every larger $\delta<1-\rho$.
\end{remark}

\section{Consequences for the grand MCA challenge}\label{sec:grand}

\begin{corollary}[universal field-size cap for the challenge envelope]\label{cor:grand}
Let $\rho\in\{\frac12,\frac14,\frac18,\frac1{16}\}$ and set
\[
N_\rho:=1024\quad(\rho\in\{\tfrac12,\tfrac14,\tfrac18\}),\qquad N_{1/16}:=2048 .
\]
Let $\F$ be any finite field with $q<2^{256}$, let $\B\subseteq\F$ be any subfield, let $D\subseteq\B^\times$ be a multiplicative coset of order $n$ with $N_\rho\mid n$, and let $k=\rho n\le2^{40}$. Then $C=\RS[\F,D,k]$ satisfies
\[
\emca(C,\delta)
>
\frac1{2k}\Bigl(1-\frac nq\Bigr)
\ge
2^{-86}
\gg
2^{-128}
\]
for every
\[
        \delta\in
        \Bigl[1-\rho-\frac2{N_\rho},\,1-\rho\Bigr).
\]
Moreover, the same lower bound holds for $\eca(C,\delta)$ throughout the integer-admissible subinterval
\[
        \delta\in
        \Bigl[1-\rho-\frac2{N_\rho},\,1-\rho-\frac1n\Bigr].
\]
If $q\ge2n$, the lower bound improves to $\ge2^{-42}$. In particular,
\[
\delta^*_C(2^{-128})\;\le\;1-\rho-2^{-9}\quad(\rho\in\{\tfrac12,\tfrac14,\tfrac18\}),
\qquad
\delta^*_C(2^{-128})\;\le\;1-\rho-2^{-10}\quad(\rho=\tfrac1{16}).
\]
\end{corollary}

\begin{proof}
We verify the hypotheses of \cref{thm:main} with $N=N_\rho$. Divisibility: $a=n/N_\rho$ divides $k$ because $k/a=\rho N_\rho\in\mathbb Z$ (equal to $512,256,128,128$ respectively); $(1-\rho)N_\rho\ge128\ge3$. For \eqref{eq:hyp}, the entropy bound $\binom N\ell\ge2^{N\Hb(\ell/N)}/(N+1)$ together with the conservative evaluations of $\Hb$ recorded in \cref{tab:numerics} gives $\log_2\binom{N_\rho}{\rho N_\rho+2}\ge552$ in all four cases, while
\[
|\B|\Bigl(\frac qk+1\Bigr)\;\le\;q\Bigl(\frac qk+1\Bigr)\;\le\;q^2+q\;<\;2^{512}+2^{256}\;<\;2^{513}\;\le\;2^{552},
\]
so \eqref{eq:hyp} holds with at least $39$ bits to spare (and over $500$ bits at rate $\frac12$). Since $D\subseteq\F^\times$, we have $n\le q-1$; for fixed $n$, the function $1-n/q$ is minimized at $q=n+1$, hence $1-n/q\ge1/(n+1)$. Also $n=k/\rho\le16k\le2^{44}$. Therefore
\[
\frac1{2k}\Bigl(1-\frac nq\Bigr)\ge\frac1{2k(n+1)}>2^{-86}.
\]
If $q\ge2n$ then $1-n/q\ge1/2$, giving the sharper $1/(4k)\ge2^{-42}$. The gap is $2/N_\rho=2^{-9}$ for $\rho\in\{1/2,1/4,1/8\}$ and $2/N_\rho=2^{-10}$ for $\rho=1/16$. \Cref{thm:main} gives the MCA lower bound throughout the stated sub-capacity interval, and gives the CA lower bound on the integer-admissible subinterval ending at $1-\rho-1/n$.
\end{proof}

\begin{table}[ht]
\centering
\begin{tabular}{ccccccc}
\toprule
$\rho$ & $N_\rho$ & $\ell_2=\rho N_\rho+2$ & $\Hb(\ell_2/N_\rho)\ge$ & $\log_2\binom{N_\rho}{\ell_2}\ge$ & needed & gap $2/N_\rho$\\
\midrule
$1/2$  & $1024$ & $514$ & $0.99998$ & $1013$ & $513$ & $2^{-9}$\\
$1/4$  & $1024$ & $258$ & $0.8143$  & $823$  & $513$ & $2^{-9}$\\
$1/8$  & $1024$ & $130$ & $0.5490$  & $552$  & $513$ & $2^{-9}$\\
$1/16$ & $2048$ & $130$ & $0.34105$ & $687$  & $513$ & $2^{-10}$\\
\bottomrule
\end{tabular}
\caption{Numerics for \cref{cor:grand}: $\log_2\binom{N}{\ell}\ge N\Hb(\ell/N)-\log_2(N+1)$, rounded down. The ``needed'' column is the universal upper bound $\log_2\bigl(|\B|(q/k+1)\bigr)<513$ valid throughout the envelope $q<2^{256}$, $k\ge1$.}
\label{tab:numerics}
\end{table}

\begin{remark}[hypothesis $N_\rho\mid n$]
For the $2$-adic domains of the challenge this holds whenever $n\ge2^{10}$ (resp.\ $2^{11}$), i.e.\ for all but very small instances. For mixed-radix smooth $n$, replace $N_\rho$ by any divisor $N\mid n$ of comparable size with $\rho N\in\mathbb Z$; the entropy count in \cref{tab:numerics} is insensitive to the exact value of $N$ at fixed $\log_2 N$.
\end{remark}

\begin{remark}[two-tier answer to the grand challenge]\label{rem:twotier}
Combining \cref{cor:grand} with the grand-challenge cap of \cite[Prop.\ ``Grand-challenge cap'']{Cho26b} (gap $\frac1{64}$, prime fields up to $\approx2^{150}$, via value-set coverage) yields the current state of the negative side for smooth multiplicative domains:
\[
\delta^*_C(2^{-128})\;\le\;
\begin{cases}
1-\rho-\frac1{64}, & \F\ \text{prime},\ q\lesssim2^{150},\\[2pt]
1-\rho-2^{-9}, & \rho\in\{\frac12,\frac14,\frac18\},\ \text{every}\ \F,\ q<2^{256},\\[2pt]
1-\rho-2^{-10}, & \rho=\frac1{16},\ \text{every}\ \F,\ q<2^{256}.
\end{cases}
\]
Together with the positive results collected in \cite[Table~1]{ABF26}, this pins the grand MCA threshold between the Johnson region and the explicit capacity-edge caps above, throughout the stated challenge envelope. The second and third tiers supersede the conditional cap of \cite[Thm.\ ``Universal cap'']{Cho26b} ($\delta^*<1-\rho-2^{-11}$ under $|\F|\ge2n$), which was the previous state in this band.
\end{remark}

\begin{corollary}[deployed parameters: KoalaBear sextic]\label{cor:deployed}
Let $\B=\F_p$ with $p=2^{31}-2^{24}+1$, let $\F=\F_{p^6}$ (so $q=p^6\approx2^{185.93}$), let $D\subseteq\B^\times$ be the subgroup of order $n=2^{21}$, and let $k=2^{20}$, $\rho=\frac12$, as in \cite[\S6.3]{ABF26}. Then
\[
\emca\bigl(\RS[\F_{p^6},D,2^{20}],\,\delta\bigr)
>
\bigl(1-2^{-164}\bigr)\,2^{-21}
>
2^{-22}
\]
for every
\[
        \delta\in\Bigl[\frac12-2^{-7},\,\frac12\Bigr).
\]
Moreover, the same lower bound holds for $\eca$ throughout the integer-admissible subinterval
\[
        \delta\in
        \Bigl[\frac12-2^{-7},\,\frac12-2^{-21}\Bigr).
\]
\end{corollary}

\begin{proof}
Apply \cref{thm:main} with $N=256$, $a=2^{13}\mid k$, $(1-\rho)N=128\ge3$, gap $2/N=2^{-7}$. Hypothesis \eqref{eq:hyp}: $\log_2\binom{256}{130}\ge256\cdot\Hb(130/256)-\log_2257\ge256\cdot0.99982-8.01>247$, while $|\B|(q/k+1)<2^{31}\,(2^{166}+1)<2^{198}$. The error bound is $\frac1{2k}(1-n/q)$ with $n/q=2^{21}/p^6<2^{-164}$. \Cref{thm:main} gives the stated MCA bound on $[\frac12-2^{-7},\frac12)$ and the stated CA bound on $[\frac12-2^{-7},\frac12-2^{-21})$.
\end{proof}

\begin{corollary}[all interleaved rows of {\cite[Tables 2--3]{ABF26}}]\label{cor:rows}
In the setting of \cref{cor:deployed}, for each $s=2^j$ with $0\le j\le12$ let $C_s:=\RS[\F_{p^6},D_s,k_s]$ with $|D_s|=n_s=2^{21}/s$ and $k_s=2^{20}/s$ (rate $\frac12$, $s\,n_s=2^{21}$ as deployed). Then for every $\delta\in\bigl[\frac12-2^{-7},\,\frac12\bigr)$,
\[
\emca\bigl(C_s^{\equiv s},\delta\bigr)
>
\frac s{2^{21}}\bigl(1-2^{-164}\bigr)
\ge
2^{-21}\bigl(1-2^{-164}\bigr).
\]
Moreover, the same lower bound holds for $\eca(C_s^{\equiv s},\delta)$ on the integer-admissible subinterval
\[
        \delta\in
        \Bigl[\frac12-2^{-7},\,\frac12-\frac1{n_s}\Bigr).
\]
Every row of the deployed parameter family therefore fails the $2^{-128}$ target at gap $2^{-7}$ below capacity.
\end{corollary}

\begin{proof}
For each $s\le2^{12}$: $256\mid n_s$ (as $n_s\ge2^9$), $a_s=n_s/256=2^{13}/s$ divides $k_s=2^{20}/s$, and $(1-\rho)\cdot256=128\ge3$. Hypothesis \eqref{eq:hyp} for $C_s$: $|\B|(q/k_s+1)<2^{31}(2^{166}\,s+1)<2^{198}\,s\le2^{210}\le2^{247}\le\binom{256}{130}$. \Cref{thm:main} gives the MCA lower bound for $C_s$ throughout $[\frac12-2^{-7},\frac12)$, and the CA lower bound throughout $[\frac12-2^{-7},\frac12-1/n_s]$. \Cref{lem:inter} transfers the corresponding lower bounds to the $s$-interleaved code.
\end{proof}

\begin{proposition}[independent slacked variant via \cref{thm:B}]\label{prop:slacked}
In the setting of \cref{thm:main}, assume instead $(1-\rho)N\ge2$, $n\ge3N$, and
\[
\binom{N}{\rho N+1}\;\ge\;|\B|\cdot q .
\]
Then
\[
\eca\Bigl(C,\ 1-\rho-\frac1N+\frac2n,\ 1-\rho-\frac1n\Bigr)\;\ge\;\frac1{2n}.
\]
In particular, with $N=N_\rho$ as in \cref{cor:grand} and $n\ge3N_\rho$, every challenge-envelope instantiation satisfies this slacked CA lower bound, and $\frac1{2n}\ge2^{-45}\gg2^{-128}$.
\end{proposition}

\begin{proof}
Set $\delta_0:=1-\rho-\frac1N$. By $(1-\rho)N\ge2$, $\delta_0\in(0,1-\rho)$. By $n\ge3N$, we also have $\delta_0+2/n\le1-\rho-1/n$, so the parameters in \cref{thm:B} are admissible. By \cref{lem:fiber}(i), some $u_z$ carries at least $\binom N{\rho N+1}/|\B|\ge q$ codewords of $C$ within radius $\delta_0$, so $\Lst(C,\delta_0)\ge q$. The contrapositive of \cref{thm:B} gives
\[
\eca\Bigl(C,\delta_0+\frac2n,1-\rho-\frac1n\Bigr)\ge\frac1{2n}.
\]
For the numerical claim, the worst case of $\log_2\binom{N_\rho}{\rho N_\rho+1}$ over the four rates is at $\rho=1/8$, where the entropy bound used in \cref{tab:numerics} certifies $\log_2\binom{1024}{129}\ge1024\,\Hb(129/1024)-\log_2 1025\ge549$ (the exact value is $\approx554.7$); hence $\binom{N_\rho}{\rho N_\rho+1}\ge2^{549}>2^{512}>|\B|q$ throughout $q<2^{256}$. Finally $n=k/\rho\le16k\le2^{44}$, so $1/(2n)\ge2^{-45}$.
\end{proof}

\begin{remark}
\Cref{prop:slacked} is logically independent of \cref{thm:A} and yields a weaker error ($\frac1{2n}$ versus $\frac1{2k}$) in the \emph{maximally slacked} form of correlated agreement --- the internal radius sits at capacity minus $\frac1n$ --- which is the weakest failure mode one can certify. That even this fails by $83$ orders of magnitude (in bits) above the $2^{-128}$ target gives a separate fallback route through the imported BCHKS/ABF theorem.
\end{remark}

\section{Subfield confinement and the shape of certifying lines}\label{sec:subfield}

The deployed instantiations run the protocol over an extension $\F/\B$ while keeping the domain in the base: $D\subseteq\B$. The subfield-confinement theorem of \cite[Thm.\ ``Subfield confinement'']{Cho26b} pins down where bad slopes can live for lines whose words are base-rational: all support-wise MCA-bad slopes lie in $\B$. The following lemma refines this by tracking two radii, so that it also covers the proximity-loss form of correlated agreement (\cref{def:ca}); part (b) is the statement of \cite{Cho26b}, reproved here for completeness in the present notation.

\begin{lemma}[subfield confinement; refines {\cite{Cho26b}}]\label{lem:confine}
Let $\B\subseteq\F$ be finite fields, $D\subseteq\B^\times$, $C=\RS[\F,D,k]$, and let $f,g\in\B^n\subseteq\F^n$ be $\B$-valued words. Then:
\begin{enumerate}[label=\textup{(\alph*)}]
\item for any $0\le\delta_{\mathrm{fld}}\le\delta_{\mathrm{int}}<1$, every CA-bad slope for $(f,g)$ at radii $(\delta_{\mathrm{fld}},\delta_{\mathrm{int}})$ lies in $\B$;
\item for any $\delta\in(0,1)$, every MCA-bad slope for $(f,g)$ at radius $\delta$ lies in $\B$.
\end{enumerate}
Consequently the bad-slope set of any $\B$-valued pair has density at most $|\B|/q$ in $\F$. The same holds for pairs $(\lambda f,\lambda g)$ with $\lambda\in\F^\times$ and $f,g$ $\B$-valued, by $\F$-linearity of $C$.
\end{lemma}

\begin{proof}
Fix a $\B$-basis $1=\beta_0,\beta_1,\dots,\beta_{m-1}$ of $\F$. Every word $w\in\F^n$ decomposes uniquely as $w=\sum_i\beta_i w_i$ with $w_i\in\B^n$, and every codeword $c\in\RS[\F,D,k]$ decomposes as $c=\sum_i\beta_i c_i$ with $c_i\in\RS[\B,D,k]$: write the coefficient vector of the underlying polynomial in the basis and evaluate at $D\subseteq\B$.

Let $z\in\F\setminus\B$, say $z=\sum_iz_i\beta_i$ with $z_{i^*}\ne0$ for some $i^*\ge1$. We first record the support-wise consequence. Suppose that $c\in C$ agrees with $f+zg$ on some set $S\subseteq D$. Since $f_j,g_j\in\B$, the basis components of $(f+zg)_j$ are $f_j+z_0g_j$ (component $0$) and $z_ig_j$ (components $i\ge1$). Agreement on $S$ holds componentwise, so on $S$:
\[
f+z_0g=c_0,\qquad z_{i^*}\,g=c_{i^*}.
\]
Hence $g$ agrees on $S$ with $z_{i^*}^{-1}c_{i^*}\in C$, and then $f=(f+z_0g)-z_0g$ agrees on $S$ with $c_0-z_0z_{i^*}^{-1}c_{i^*}\in C$. Thus every support on which the line point $f+zg$ is explained by $C$ also explains the pair $(f,g)$ on that same support.

For CA, if $\dist(f+zg,C)\le\delta_{\mathrm{fld}}$, choose such an $S$ with $|S|\ge(1-\delta_{\mathrm{fld}})n$; the preceding paragraph gives $\dist_2((f,g),C^{\equiv2})\le\delta_{\mathrm{fld}}\le\delta_{\mathrm{int}}$, so $z$ is not CA-bad. For MCA, the preceding paragraph rules out every possible bad support $S$ in the support-wise definition. Hence $z$ is not MCA-bad either.
\end{proof}

\begin{corollary}[certifying lines over extensions are genuinely $\F$-valued]\label{cor:Fvalued}
At the parameters of \cref{cor:deployed}, $|\B|/q=p^{-5}<2^{-154}$. Hence any pair $(f,g)$ whose CA- or MCA-bad-slope density at some radius
\[
        \delta\in
        \Bigl[\frac12-2^{-7},\,\frac12\Bigr)
\]
exceeds $2^{-154}$ contains a word that is not $\B$-valued, even after removing a common scalar factor. In particular, any pair attaining the MCA lower bound $>2^{-22}$ of \cref{cor:deployed} is genuinely $\F$-valued. On the integer-admissible subinterval
\[
        \delta\in
        \Bigl[\frac12-2^{-7},\,\frac12-2^{-21}\Bigr),
\]
the same statement applies to pairs attaining the CA lower bound. All explicit failure witnesses currently known on these domains (the locator lines of \cite{Cho26a,Cho26b} and the fiber words of \cref{lem:fiber}) are $\B$-rational and are therefore inert here: the failure certified by \cref{cor:deployed} is fiber-borne and, at present, non-constructive.
\end{corollary}

\begin{proof}
Immediate from \cref{lem:confine} and \cref{cor:deployed}: a $\B$-valued pair, up to common scalar, has CA- and MCA-bad-slope density at most $|\B|/q=p^{-5}<2^{-154}<2^{-22}$. Thus any pair witnessing the deployed lower bound cannot be of this form.
\end{proof}

\begin{remark}
\Cref{lem:confine,cor:Fvalued} sharpen the division of labor between the two grand challenges of \cite{ABF26} over deployed extensions: the \emph{list} side does not deflate ($\RS[\B,D,k]\subseteq\RS[\F,D,k]$, so all base-field list lower bounds, including \cref{lem:fiber}, persist verbatim), while every sub-$2^{-128}$ \emph{MCA} attack by explicit base-rational witnesses is killed by confinement --- and yet MCA still fails at $2^{-22}$.  In earlier drafts this failure was only extracted non-constructively from \cref{thm:A}; \cref{cor:extension-pole-deep-list-floor,cor:extension-pole-quotient-remainder-floor} now give a constructive source whenever a large list is supplied: choose the simple pole in $\F\setminus\B$ and print its value-bucket table. The deflation and decoupling corollaries of \cite{Cho26b} record the same division at the no-loss radius, using the same corrected rounding $p^{-5}=2^{-154.94\ldots}<2^{-154}$ as in \cref{cor:Fvalued}. Exhibiting deployed-size bucket tables for such lines is Open Problem~\ref{prob:explicit}.
\end{remark}


\section{Prize-facing status and a completion program}\label{sec:prize-program}

The Proximity Prize asks for the largest admissible radius, not only for a
near-capacity obstruction.  Thus the present paper should be read as a
one-sided theorem: it proves a universal unsafe band, while the remaining work
is to prove a matching safe band below that cap, or to find additional
obstruction floors that push the threshold even farther from capacity.

\subsection{What is already settled by this paper}

For the grand MCA challenge at target error $\varepsilon^*=2^{-128}$, the
unconditional output of \cref{cor:grand} is the following negative certificate:
for every challenge-envelope row satisfying the printed divisor hypothesis,
\[
        \delta\ge 1-\rho-2^{-9}
        \quad(\rho\in\{1/2,1/4,1/8\})
\]
and
\[
        \delta\ge 1-\rho-2^{-10}
        \quad(\rho=1/16)
\]
are impossible MCA radii.  Equivalently, any full determination of
$\delta_C^*(2^{-128})$ for these smooth multiplicative rows must search below
these gaps.  This conclusion uses only the self-contained deep-point conversion,
the quotient-fiber pigeonhole, and support-wise MCA monotonicity; it no longer
rests on importing the Crites--Stewart conversion.

The theorem also distinguishes two different grades of information.  First, the
\emph{threshold grade} is already enough for the prize inequality: the certified
error is $\gg2^{-128}$.  Second, the \emph{failure-profile grade} is still weak:
the present proof gives error of order $1/k$, not error $1-o(1)$.  Closing the
error-one profile at fixed gaps such as $1/64$ is a different, stronger problem.
It would improve our understanding of the failure landscape, but is not needed
to obtain the stated upper bound on $\delta_C^*(2^{-128})$.

\subsection{Integer-staircase formulation}

Let
\[
        a(\delta):=\lceil(1-\delta)n\rceil,
        \qquad r(\delta):=n-a(\delta)
\]
be the agreement and integer-radius forms of the same radius.  For a fixed finite-slope line
field $F_{\rm line}$, write $q_{\rm line}=|F_{\rm line}|$ and let
\[
        B_{\rm mca}(a)
\]
be the exact maximum number of MCA-bad line parameters at agreement at least
$a$, for the line family used in the challenge.  Then
\[
        \emca(C,1-a/n)=B_{\rm mca}(a)/q_{\rm line}
\]
under the closed finite-slope convention.  For projective-slope or curve samplers, the same staircase formulation applies after replacing $q_{\rm line}$ and $B_{\rm mca}$ by the corresponding sampler denominator and numerator.  Thus the MCA prize problem is an
integer staircase problem: find the first agreement $a$ for which
\[
        B_{\rm mca}(a)\le 2^{-128}q_{\rm line}.
\]
The present paper supplies lower bounds on $B_{\rm mca}(a)$ near capacity by
routing quotient-fiber list mass through deep points.

\begin{proposition}[one-step threshold certificate]\label{prop:onestep}
Fix a row $C$ and a denominator $q_{\rm line}$.  Suppose that for some agreement
$a_0$ one has a lower certificate
\[
        B_{\rm mca}(a_0)>2^{-128}q_{\rm line}
\]
and an upper theorem
\[
        B_{\rm mca}(a_0+1)\le2^{-128}q_{\rm line}.
\]
Then, under the closed integer-radius convention, radius $1-a_0/n$ is unsafe and
radius $1-(a_0+1)/n$ is safe.  As a real-radius supremum, the threshold is
$1-a_0/n$ and is not attained if the closed ball at that endpoint is counted.
\end{proposition}

\begin{proof}
The function $B_{\rm mca}(a)$ is nonincreasing in $a$, since raising the
required agreement only removes witnesses.  At the closed radius $1-a/n$, the
MCA error is $B_{\rm mca}(a)/q_{\rm line}$.  The two displayed inequalities are
therefore exactly the unsafe/safe transition at consecutive integer agreement
levels.
\end{proof}

\subsection{Conditional MCA completion theorem}

The following theorem is a deliberately explicit proof blueprint.  Each
hypothesis is a concrete mathematical task, not a black-box ``large field''
assumption.

\begin{theorem}[conditional MCA threshold theorem]\label{thm:conditional-mca}
Fix a smooth multiplicative row $C=\RS[F,D,k]$, a line field $F_{\rm line}$, and
a target $\varepsilon^*=2^{-128}$.  Suppose the following four inputs are proved
for this row, uniformly in the received line:
\begin{enumerate}[label=(\roman*),leftmargin=2.2em]
\item \textbf{Quotient-profile lower and upper ledger.}  For every divisor
$c\mid n$ and every agreement $a=mc+s$ with $0\le s<c$, the quotient-periodic
bad slopes, including remainder supports consisting of $m$ full fibers plus
$s$ residual points, are counted with overlaps across divisors accounted for.
Write the resulting divisor-lattice upper numerator as $B_Q^{\rm all}(a)$;
individual divisor contributions $B_Q(a;c)$ may also be printed, but the upper
ledger must count their union or a certified safe sum rather than only their
maximum.
\item \textbf{Aperiodic residue-line packing.}  After removing tangent and
quotient-periodic locators, the Hankel/residue-line incidence variety has at
most $B_{\rm ap}(a)$ bad line parameters.
\item \textbf{Extension-line accounting.}  If $F_{\rm line}$ is a proper
extension of the field generated by $D$, either all genuinely extension-valued
line witnesses are included in $B_{\rm ap}(a)$, or a separate extension term
$B_{\rm ext}(a)$ is proved.
\item \textbf{No missing line family.}  The line family in the theorem is the
same finite-slope or projective-slope family used by the challenge or protocol
under consideration.
\end{enumerate}
Assume moreover that the resulting upper numerator
\[
        B^+(a):=B_{\rm tan}(a)+B_Q^{\rm all}(a)
                +B_{\rm ap}(a)+B_{\rm ext}(a)
\]
and a corresponding lower numerator $B^-(a)$ satisfy
\[
        B^-(a_0)>\varepsilon^* q_{\rm line},
        \qquad
        B^+(a_0+1)\le\varepsilon^* q_{\rm line}
\]
for some $a_0$.  Then $\delta_C^*(\varepsilon^*)$ is determined at the
one-step accuracy of \cref{prop:onestep}.
\end{theorem}

\begin{proof}
The four hypotheses classify every possible MCA-bad line parameter into one of
four buckets: tangent, quotient-periodic, aperiodic, or genuinely extension-line.
The quotient bucket is counted by the divisor-lattice union ledger $B_Q^{\rm all}$,
so overlaps or simultaneous quotient descriptions cannot invalidate the upper
bound.  The displayed $B^+(a)$ is therefore an upper bound for $B_{\rm mca}(a)$, while
$B^-(a)$ is a certified lower bound.  The conclusion is exactly
\cref{prop:onestep}.
\end{proof}

The present paper supplies part of the lower numerator $B^-(a)$ in the
near-capacity band: quotient-fiber list mass forces MCA mass above the printed
universal caps.  It does not prove the aperiodic upper input (ii), and therefore
does not by itself determine the exact staircase.

\subsection{Conditional interleaved-list completion theorem}

The grand list-decoding challenge asks for the largest $\delta$ such that an
interleaved list, divided by the field used for the relevant challenge or list
ledger, is below $2^{-128}$.  In the simplest same-field form this is
\[
        |\Lambda(C^{\equiv m},\delta)|\le 2^{-128}|F| .
\]
The quotient-fiber construction in this paper already gives lower floors for
ordinary lists, and ordinary lists embed into interleaved lists by placing the
same received word in one row and zero words in the others.  On the MCA side,
\cref{sec:aperiodic-hankel-certificates} supplies a certified eliminant route
for aperiodic Hankel packing, but it does not prove that the needed eliminants
always exist.  The missing positive list direction is an arbitrary-word
locator-fiber theorem.

\begin{theorem}[conditional interleaved-list threshold theorem]\label{thm:conditional-list}
Fix an interleaving arity $m$.  Suppose that for every received word $U$ and
agreement $a=k+\sigma$ one has a base-code list bound of the form
\[
        |\Lambda(C,1-a/n,U)|
        \le
        L_Q(a,U)+L_{\rm ap}(a),
\]
where $L_Q$ is an explicit quotient-profile list contribution and
$L_{\rm ap}$ is a uniform aperiodic locator-fiber bound.  Suppose also that the
same bound is available at the higher agreement levels that occur in the
codegree recursion, starting with $2a-k$.  Then the interleaved-list threshold is
obtained by applying the codegree recursion row by row and comparing the
resulting numerator to $2^{-128}|F|$.
\end{theorem}

\begin{proof}
For two rows, split the second-row list according to whether its agreement set
has size below or above $2a-k$.  Below $2a-k$, two first-row codewords would
agree on more than $k$ positions and hence are equal; above $2a-k$, use the
base-code list bound at agreement $2a-k$.  Iterating gives the stated $m$-row
recursion.  This is the same codegree mechanism used in the experimental ledger;
the only missing input is the arbitrary-word locator-fiber bound.
\end{proof}

\subsection{Concrete proof obligations}

The fastest route from the present cap to a full prize solution is to turn the
following proof obligations into theorems.  The lower-floor half of T1 is now supplied in scanner-ready form by
\cref{lem:heaviest-prefix-locator-floor,thm:quotient-remainder-deep-floor}: every
quotient-plus-remainder agreement $a=mc+s$ has a generated-field prefix-fiber
floor, an image-size fallback, an ambient-box fallback, and a converted CA/MCA
lower numerator.  The support-family upper half is exact at the divisor-lattice
level by \cref{thm:exact-divisor-block-support-ledger}: overlaps between quotient
descriptions at different divisor scales are counted by a common-block
coefficient extraction rather than by a raw safe sum.  The distinct-parameter
branch image is exact inside every declared quotient family by the lcm ledger of
\cref{thm:exact-quotient-image-lcm-ledger,cor:affine-line-lcm-ledger,cor:projective-line-lcm-ledger}; duplicate finite slopes, projective slopes, and
finite curve parameters are coalesced before the numerator is printed.  T2 is now
also partially upgraded on its regular overdetermined branch by
\cref{thm:canonical-affine-rankdrop-ledger,cor:canonical-regular-closed-ball-hankel-packing,thm:canonical-curve-rankdrop-ledger}: instead of choosing one regular
Hankel minor, the scanner takes the gcd of all maximal minors, then an lcm over
exact agreements, and finally removes exact quotient/tangent paid roots.  What
remains open is the genuinely structural safe-side theorem: proving that every
bad parameter in the intended near-capacity regime either belongs to the tangent
or quotient image, or is bounded by the pivot-chart aperiodic and extension
ledgers.  The underdetermined and rank-drop singular Hankel buckets remain the
named obstruction for T2.  T5 is partially upgraded as well: finite-parameter
curve and projective-line quotient branches have exact gcd/lcm image ledgers,
and finite-parameter curve regular rank-drop branches have canonical maximal-minor
gcd ledgers, while the aperiodic curve and projective pivot charts still require
eliminant certificates.

\begin{enumerate}[label=\textbf{T\arabic*.},leftmargin=2.5em]
\item \textbf{Remainder quotient-profile theorem.}  Generalize the fiber lemma
from full quotient-fiber unions to all agreements $a=mc+s$.  The target is an
exact lower and upper quotient ledger over the divisor lattice of $D$.  Version
10 supplies the scanner-ready lower half via heaviest prefix fibers and the
support-union upper half via the exact divisor-block coefficient formula; inside
a declared quotient family, the gcd/lcm image ledger also coalesces duplicate
finite, projective, and curve parameters.  What remains is the structural
exhaustion theorem saying that the declared quotient family captures all
quotient-periodic bad parameters in the intended regime, with the subfield
denominator $|B|$ printed separately from the line denominator $q$.

\item \textbf{Aperiodic Hankel-packing theorem.}  Use the syndrome/Hankel normal
form of \cite{Cho26b} to prove that, after quotient-periodic and tangent
locators are removed, the number of remaining bad line parameters is at most
$B_{\rm ap}(a)$.  The regular overdetermined nonsingular branch is now
canonical by \cref{thm:canonical-affine-rankdrop-ledger,cor:canonical-regular-closed-ball-hankel-packing}: a scanner takes the gcd of all nonzero maximal
Hankel minors and then an lcm over exact agreement levels, rather than choosing
one minor.  The remaining proof problem is concentrated in the underdetermined
and rank-drop singular buckets: construct small pivot-chart eliminants for them,
or classify their residual components as quotient, tangent, extension-field
exceptional, or genuinely new.

\item \textbf{Extension-line lift or counterexample.}  For $B\subset F$ and
$D\subset B$, either prove that genuinely $F$-valued lines are covered by the
aperiodic ledger with a stable numerator, or exhibit explicit extension-valued
counterexamples.  \Cref{cor:Fvalued} shows that deployed failures certified by
this paper must be genuinely extension-valued; \cref{cor:extension-pole-deep-list-floor,cor:extension-pole-quotient-remainder-floor}
now provide an explicit lower-side mechanism from any large list by taking the
pole in $F\setminus B$.  The remaining safe-side problem is to classify all such
extension-valued lines or charge them to the aperiodic ledger.

\item \textbf{Arbitrary-word locator-fiber theorem.}  Prove the base-code list
bound needed by \cref{thm:conditional-list}, first for monomial-prefix received
words and quotient-free dimensions, then for arbitrary received words with
quotient cores budgeted.

\item \textbf{Curve and projective-line ledgers.}  Extend the affine-line Hankel
classification to projective slopes and to finite-parameter power curves
$f_0+\gamma f_1+\cdots+\gamma^d f_d$.  Version 10 gives exact gcd/lcm ledgers
for projective and finite-curve quotient branches and canonical rank-drop gcds
for regular overdetermined finite-curve buckets.  The remaining work is the
aperiodic curve and projective pivot atlas, including exceptional
vanishing-vector strata.

\item \textbf{Certificate scanner.}  Implement a finite staircase scanner that,
for each challenge row, prints tangent, quotient, generated-field entropy,
aperiodic, extension, curve, and interleaved-list numerators at every agreement
$a$.  The desired output is an interval
\[
        a_{\rm unsafe}<a_C^*\le a_{\rm safe}
\]
that shrinks to one step when \cref{thm:conditional-mca} or
\cref{thm:conditional-list} applies.

\item \textbf{Formal finite core.}  Formalize the self-contained deep-point
conversion, the quotient-fiber pigeonhole, support-wise MCA monotonicity, and
the integer staircase lemma.  These are small enough for proof-assistant
checking and directly strengthen a prize submission.
\end{enumerate}

\subsection{Recommended submission posture}

As a standalone prize-facing submission, this paper should claim a significant
partial result: a universal negative band for the grand MCA challenge over
smooth multiplicative rows in the stated field-size envelope.  It should not
claim the full prize threshold, a positive theorem below the cap, protocol
soundness, or a curve-MCA theorem.  The full prize claim should be made only
after the conditional ledgers above are replaced by theorems and the staircase
scanner produces a one-step unsafe/safe transition for the relevant rows.


\section{Discussion}\label{sec:discussion}

\paragraph{Values versus fibers.}
Above the per-rate value-set reach of \cite{Cho26b} (roughly $2^{150}$--$2^{162}$), the value route --- showing that the slope multiset $\eone(\ell^\wedge Q)$ covers a $2^{-128}$ fraction of $\F$ at a fixed large prime --- runs into the per-fiber collision problem (\cite[Rem.\ ``The remaining band'']{Cho26b}): characteristic-zero locator values can collide arbitrarily badly under reduction at a worst-case prime. \cite{Cho26b} therefore closed the band at the $\delta^*$ grade only conditionally and with weaker constants --- error $2^{-42}$ at gap $2^{-11}$, assuming $|\F|\ge2n$, via quotient-core lists and the same conversion \cref{thm:A} --- leaving the \emph{error-one} grade open. The present route also never counts values, and improves each parameter of that cap: pigeonhole guarantees one heavy fiber regardless of how the values $\eone(A)$ distribute, \cref{thm:A} converts heavy fibers into correlated-agreement error directly, and the subfield pigeonhole carries the construction to extension fields with no condition relating $n$ and $|\F|$. The cost relative to the error-one results of \cite{Cho26a,Cho26b} below $2^{150}$ is quantitative: error $\approx\frac1{2k}$ instead of $1-o(1)$, and gap $2^{-9}$ (uniformly $2^{-10}$ across all four rates) instead of $\frac1{64}$. For the prize question --- is $\emca\le2^{-128}$ achievable above the stated gaps? --- this cost is immaterial.

\paragraph{Limits of the method.}
Hypothesis \eqref{eq:hyp} needs $N\,\Hb(\rho)\gtrsim2\log_2q-\log_2k$, so the smallest certifiable gap is
\[
\frac2N\;\approx\;\frac{2\,\Hb(\rho)}{2\log_2q-\log_2k}\;\approx\;\frac{\Hb(\rho)}{\log_2 q},
\]
and the certified error is capped at $\approx\frac1k$ by the shape of \cref{thm:A} (any $\eta<1$). Since $q$ enters only as the size of the field from which the line challenge $\gamma$ is drawn, this is also a protocol-facing floor: enlarging the challenge field cannot push a deep-point-conversion failure below gap $\approx\Hb(\rho)/\log_2 q_{\mathrm{line}}$, which quantitatively bounds the ``larger challenge field'' fallback in certificate frameworks. The construction also needs $N\le n$, i.e.\ $\log_2q\lesssim\Hb(\rho)\,n/2$; this covers the entire cryptographic regime $\mathrm{poly}(n)\le q\le2^{o(n)}$, but for $q\ge2^{\Hb(\rho)n/2}$ the method is vacuous and the value-route caps of \cite{Cho26a,Cho26b} are the only ones known. Within the band, the \emph{error-one} question at gap $\frac1{64}$ (Open Problem~\ref{prob:errorone}) remains exactly as open as before: the present result bounds $\delta^*$, not the failure profile near capacity.

\paragraph{Comparison with the cyclotomic sieve.}
The sieve of \cite{Cho26a} produces, for infinitely many primes, failures of error $(\log p)^{-6}$ at gap $\Theta(1/\log p)$, but its prime set is given by Siegel--Walfisz and is ineffective. \Cref{cor:grand} trades error strength for universality: every field below $2^{256}$, fully effective, and elementary given \cref{thm:A}. The two results triangulate the failure landscape from different directions.

\paragraph{Pressure on the packing conjecture.}
By the exact normal form of \cite{Cho26b}, $\emca(C,\delta)=\max_t\Lambda^{\mathrm{NC}}_t(\delta)/q$, where $\Lambda^{\mathrm{NC}}$ is a noncontained residue-line packing number. \Cref{cor:deployed} therefore forces $\max_t\Lambda^{\mathrm{NC}}\ge(q-n)/2k\approx2^{164}$ at gap $2^{-7}$ on the deployed sextic --- superpolynomially many noncontained residue lines at $n=2^{21}$. This puts no direct pressure on the final MCA conjecture of \cite[Conj.\ ``Final floor- and quotient-corrected MCA asymptotic'']{Cho26b}: at gap $2^{-7}$ both of its corrected-reserve hypotheses fail ($\sigma=2^{14}$ lies below $Cn/\log_2n\approx2^{16.6}$, and $\sigma\log_2 q_{\mathrm{gen}}=2^{14}\cdot31\approx2^{19}$ is far below $\log_2\binom n{k+\sigma}\approx2^{21}$), so the conjectured ceiling of the form $\bigl(n^{1+o(1)}+2^{(\beta/\Hb)\,Q_{\Hb}(\eta)}\bigr)/q$ is simply not asserted there --- the packing mass extracted through \cref{thm:A} lives in the region the framework already marks as below-floor. The genuine structural question this note raises for the framework of \cite{Cho26b} is whether conversion-extracted fiber mass persists \emph{above} the corrected reserve: whether, at reserves where the conjecture does apply, the fiber route still produces packings beyond the $n^{1+o(1)}$ aperiodic term plus the quotient-periodic term, which would force an additional fiber-borne term into the conjectured ceiling. This is open.

\paragraph{Verification caveat.}
The main cap no longer depends on importing the Crites--Stewart theorem: \cref{thm:A} proves the needed deep-point list-to-CA conversion directly, including the integer-radius condition $f\le n-k-1$.  \Cref{thm:B} remains a separate slacked fallback imported through the BCHKS/ABF route, and readers using that fallback should consult the underlying source directly.  The main composition is deliberately shallow: it uses the self-contained conversion, the locator-fiber pigeonhole, and support-wise MCA monotonicity.

\paragraph{Open problems.}

\begin{problem}\label{prob:explicit}
Exhibit explicit pairs $(f,g)$ over $\F_{p^6}$ (KoalaBear sextic, parameters of \cref{cor:deployed}) with MCA-bad-slope density $>2^{-22}$ at gap $2^{-7}$. On the integer-admissible subinterval, ask for the corresponding CA-bad-slope density as well. By \cref{cor:Fvalued} such pairs exist and are not $\B$-rational; the natural candidates are residue-line normal forms of \cite{Cho26b} with denominator $E\in\F[X]$ not defined over $\B$.
\end{problem}

\begin{problem}\label{prob:errorone}
Error one in the band: does $\emca(C,1-\rho-\frac1{64})=1-o(1)$ hold for all primes $2^{150}\le p\le2^{256}$? This is the per-fiber collision problem proper; \cref{cor:grand} caps $\delta^*$ there but says nothing above error $\approx\frac1{2k}$.
\end{problem}

\begin{problem}
Improve the threshold of \cref{thm:A}: any strengthening of the hypothesis range $\eca\le\eta/k$ toward $\eca\le\eta\cdot\mathrm{polylog}(q)/k$, or of the list conclusion, immediately tightens the certified error of \cref{thm:main} toward $1/\mathrm{poly}\log$; the gap, by contrast, is governed by the binomial in \eqref{eq:hyp} and would need a denser fiber construction to shrink.
\end{problem}

\begin{problem}
Beyond multiplicative smoothness: \cref{lem:fiber} uses only that $D$ is a union of complete fibers of a low-degree map with lacunary power structure. Carry the composition to isogeny-smooth domains --- in particular the circle-group domains of Circle STARKs over $p=2^{31}-1$, where Chebyshev locators $\prod_b(T_{2^j}(x)-x_b)$ play the role of $\prod_b(X^a-b)$ --- to obtain the same universal cap for that deployed family.
\end{problem}

\paragraph{Acknowledgements.} This note builds on the Crites--Stewart conversion perspective \cite{CS25}, the BCHKS theorem \cite{BCHKS25}, and the locator framework of \cite{Cho26a,Cho26b}, following the synthesis of Arnon--Boneh--Fenzi \cite{ABF26}; it proves the main conversion needed here directly, sharpens the conditional universal-cap theorem, and refines the subfield-confinement theorem of \cite{Cho26b}.


\begin{thebibliography}{BCHKS25}

\bibitem[PP26]{ProximityPrize}
The Proximity Prize.
\newblock The Proximity Prize: \$1M Reed--Solomon Challenge.
\newblock Preliminary version, 2026.
\newblock \url{https://proximityprize.org/}.


\bibitem[ABF26]{ABF26}
G.~Arnon, D.~Boneh, and G.~Fenzi.
\newblock Open problems in list decoding and correlated agreement.
\newblock IACR Cryptology ePrint Archive, Report 2026/680, 2026.
\newblock \url{https://eprint.iacr.org/2026/680}.

\bibitem[BCHKS25]{BCHKS25}
E.~Ben-Sasson, D.~Carmon, U.~Hab{\"o}ck, S.~Kopparty, and S.~Saraf.
\newblock On proximity gaps for Reed--Solomon codes.
\newblock IACR Cryptology ePrint Archive, Report 2025/2055, 2025.
\newblock \url{https://eprint.iacr.org/2025/2055}.

\bibitem[Cho26a]{Cho26a}
P.~Chojecki.
\newblock Capacity-edge obstructions to Reed--Solomon mutual correlated agreement over smooth multiplicative domains.
\newblock Manuscript, 2026.

\bibitem[Cho26b]{Cho26b}
P.~Chojecki.
\newblock Slack, quotient cores, and the entropy gap for smooth-domain Reed--Solomon codes: list fibers, cyclotomic rigidity, Galois-amplified collisions, and the corrected slack mutual-correlated-agreement theory.
\newblock Manuscript, merged corrected edition, 2026.

\bibitem[CS25]{CS25}
E.~Crites and A.~Stewart.
\newblock On Reed--Solomon proximity gaps conjectures.
\newblock IACR Cryptology ePrint Archive, Report 2025/2046, 2025.
\newblock \url{https://eprint.iacr.org/2025/2046}.

\end{thebibliography}

\end{document}
